\documentclass[11pt,letterpaper]{article}
\usepackage[T1]{fontenc}
\usepackage[utf8]{inputenc}
\usepackage{lmodern,microtype}
\usepackage[margin=1in,headheight=15pt]{geometry}
\usepackage{amsmath,amssymb,amsthm,mathtools,bm,mathrsfs}
\usepackage{booktabs,longtable,array,enumitem}
\usepackage[dvipsnames]{xcolor}
\usepackage{fancyhdr,aliascnt}
\usepackage{hyperref}
\usepackage[nameinlink,noabbrev]{cleveref}
\hypersetup{colorlinks=true,linkcolor=MidnightBlue,citecolor=MidnightBlue,urlcolor=MidnightBlue,
 pdftitle={Surcomplex Analysis: Local Power Series, Coherent Hahn Families, and Infinitesimal Zero Geometry},
 pdfauthor={Research exposition prepared with ChatGPT}}
\pagestyle{fancy}\fancyhf{}
\fancyhead[L]{\small\scshape Surcomplex analysis}
\fancyhead[R]{\small Local series and global coherence}
\fancyfoot[C]{\thepage}
\setlength{\parskip}{4pt}
\setlength{\parindent}{1.2em}
\setlist{itemsep=2pt,topsep=5pt}
\allowdisplaybreaks[2]
\numberwithin{equation}{section}
\newtheorem{theorem}{Theorem}[section]
\newaliascnt{proposition}{theorem}
\newtheorem{proposition}[proposition]{Proposition}
\aliascntresetthe{proposition}
\crefname{proposition}{proposition}{propositions}
\newaliascnt{lemma}{theorem}
\newtheorem{lemma}[lemma]{Lemma}
\aliascntresetthe{lemma}
\crefname{lemma}{lemma}{lemmas}
\newaliascnt{corollary}{theorem}
\newtheorem{corollary}[corollary]{Corollary}
\aliascntresetthe{corollary}
\crefname{corollary}{corollary}{corollaries}
\theoremstyle{definition}
\newaliascnt{definition}{theorem}
\newtheorem{definition}[definition]{Definition}
\aliascntresetthe{definition}
\crefname{definition}{definition}{definitions}
\newaliascnt{example}{theorem}
\newtheorem{example}[example]{Example}
\aliascntresetthe{example}
\crefname{example}{example}{examples}
\newaliascnt{construction}{theorem}
\newtheorem{construction}[construction]{Construction}
\aliascntresetthe{construction}
\crefname{construction}{construction}{constructions}
\theoremstyle{remark}
\newaliascnt{remark}{theorem}
\newtheorem{remark}[remark]{Remark}
\aliascntresetthe{remark}
\crefname{remark}{remark}{remarks}
\newcommand{\No}{\mathbf{No}}
\newcommand{\Oz}{\mathbf{Oz}}
\newcommand{\K}{\mathbb K}
\newcommand{\R}{\mathbb R}
\newcommand{\C}{\mathbb C}
\newcommand{\N}{\mathbb N}
\newcommand{\Z}{\mathbb Z}
\newcommand{\Q}{\mathbb Q}
\newcommand{\OO}{\mathcal O}
\newcommand{\HH}{\mathscr H}
\newcommand{\mm}{\mathfrak m}
\newcommand{\PP}{\mathcal P}
\newcommand{\dd}{\,\mathrm d}
\newcommand{\ev}[1]{#1^{\sharp}}
\newcommand{\tmon}[1]{t^{#1}}
\DeclareMathOperator{\supp}{supp}
\DeclareMathOperator{\st}{st}
\DeclareMathOperator{\ord}{ord}
\DeclareMathOperator{\Res}{Res}
\DeclareMathOperator{\Exp}{Exp}
\DeclareMathOperator{\Log}{Log}
\DeclareMathOperator{\Arg}{Arg}
\DeclareMathOperator{\Fin}{Fin}
\DeclareMathOperator{\PI}{PI}
\DeclareMathOperator{\id}{id}
\DeclareMathOperator{\Rea}{Re}
\DeclareMathOperator{\Ima}{Im}
\DeclareMathOperator{\Span}{span}
\newcommand{\abs}[1]{\left|#1\right|}
\newcommand{\norm}[1]{\left\|#1\right\|}
\newcommand{\co}[2]{[#1^{#2}]}
\begin{document}
\begin{titlepage}
\vspace*{0.7in}
{\noindent\color{MidnightBlue}\Large\scshape A mathematical research exposition\par}
\vspace{0.35in}
{\noindent\Huge\bfseries Surcomplex Analysis\par}
\vspace{0.18in}
{\noindent\LARGE Local Power Series, Coherent Hahn Families,\par}
\vspace{0.08in}
{\noindent\LARGE and Infinitesimal Zero Geometry\par}
\vspace{0.4in}
{\noindent\large September 21, 2026\par}
\vspace{0.35in}
\begin{abstract}
The algebraic closure $\K=\No[i]$ of Conway's surreal numbers is a natural
candidate for complex analysis at infinite and infinitesimal scales. Its
neighborhood topology, however, does not support nontrivial set-indexed
convergent sequences or paths from the ordinary real interval. Consequently,
one must specify what power-series summation, analytic continuation, and
contour integration mean.

We develop a layered theory. First, normal-form summation realizes every
formal power series with surcomplex coefficients as a local analytic germ.
We prove differentiation, the Cauchy--Riemann criterion in the real-analytic
category, inverse and implicit function theorems, a ramification normal form,
and a formal residue calculus. Second, Hahn series with ordinary holomorphic
coefficient functions supply a coherent global category on standard-part
halos and on their affine rescalings. Every local formal germ embeds in this
category after a sufficiently small change of scale. In it we prove Cauchy
and Morera theorems, a constructive Weierstrass preparation and division
theorem, infinitesimal Rouch\'e stability, exact argument and residue formulas,
and recovery of the polynomial of a zero cluster from contour moments.

Counterexamples delimit the claims: unrestricted local analyticity does not
imply a global identity theorem or Liouville's theorem; even coherent halo
functions require care with Schwarz and boundedness assertions. Established
surreal foundations and the known surcomplex exponential are explicitly
credited. The proposed organization and its proved consequences are not
presented as a verified claim of publication priority.
\end{abstract}
\vfill
\noindent\small\textbf{Reading guide.} The principal constructive result is
\cref{thm:preparation}; its exact zero-counting and contour consequences are
\cref{thm:rouche,thm:argument,thm:residue}. Foundational conventions are in
\cref{sec:foundations,sec:sums}. The word ``integral'' has the explicit
coefficientwise meaning of \cref{def:integral}, not an unstated Riemann-limit
meaning on a proper class.
\end{titlepage}
\tableofcontents
\clearpage

\section{Aim, provenance, and the three layers}
\label{sec:aim}

A surcomplex number is $x+iy$ with $x,y\in\No$ and $i^2=-1$.
The point of enlarging $\C$ in this way is not merely to have more algebraic
numbers: $\C$ is already algebraically closed. The enlargement introduces
independent orders of magnitude, arbitrary infinitesimal scales, and values
at infinite arguments. An analytic theory should use this extra structure
without importing false completeness assumptions.

There is substantial prior work. Conway and Gonshor provide the arithmetic,
normal forms, and exponential background \cite{Conway,Gonshor}. Alling
explicitly develops analytic functions of surreal and surcomplex variables
\cite{Alling}. Berarducci and Mantova formulate surreal analyticity using
summable normal forms; their Section 3 establishes realization of formal
series, differentiation, and composition, and records the obstruction to
unrestricted analytic continuation \cite{BMgerms}. These foundational ideas
are not new claims of the present article. Our local arguments give a
complex-valued version with an explicit scale-separation proof.

The global construction studied here is the algebra of expressions
\begin{equation}\label{eq:introH}
 F(Z)=\sum_{\gamma\in S}f_\gamma(Z)t^\gamma,
 \qquad t^\gamma:=\omega^{-\gamma},\qquad f_\gamma\in\OO(U),
\end{equation}
where $S\subset\No$ is a set and is well ordered in its usual order.
The coefficients vary holomorphically on an ordinary complex domain $U$;
the monomials carry the surreal scales. This is a deliberate coherence
condition, not a consequence of pointwise differentiability.

Our three layers are as follows.
\begin{enumerate}[label=\textbf{\arabic*.},leftmargin=*]
\item \textbf{Intrinsic local analysis.} A function is locally a normally
summable series with arbitrary coefficients in $\K$. The algebra of germs
at a point is exactly $\K[[X]]$.
\item \textbf{Coherent Hahn analysis.} Expressions \eqref{eq:introH} are
realized on $U^\sharp$, the points with standard part in $U$. Ordinary
holomorphic dependence of every coefficient supplies global coherence.
Affine changes of scale bring this theory to arbitrary surcomplex centers.
\item \textbf{Canonical classical lifts.} An ordinary holomorphic function
has a Taylor lift to $U^\sharp$. These special coherent functions preserve
classical composition, conformal equivalence, and the Schwarz inequality.
\end{enumerate}

The distinction is substantive. Every germ in the first layer occurs in a
sufficiently small affine chart of the second. But arbitrary germs assigned
independently to different monads need not arise from one coherent family.
Similarly, a small Hahn perturbation of a classical disk map need not satisfy
all inequalities of the original disk map.

\paragraph{What is proved and what is imported.}
We use the real-closedness and normal-form theorem for $\No$, the
Hahn--Neumann summability lemma, the Kruskal--Gonshor real exponential, and
ordinary complex analysis \cite{Ahlfors} as established inputs. All passage from these
inputs to the stated surcomplex results is proved below. In particular, the
preparation theorem is established by an explicit homogeneous recursion;
no sequential contraction principle is assumed. The computations packaged
with this article check examples and are not substitutes for the proofs.

\paragraph{Scope of novelty.}
This is a self-contained development relative to the stated foundational
inputs, not a claim to have invented surreal analyticity or the surcomplex
exponential. Its particular coherent category, scale-embedding formulation,
and package of constructive zero-cluster results are developed here, but a
comprehensive priority determination for each formulation has not been made.
Questions left as research directions are separated from the theorems.

\section{The field, its valuation, and the topology}
\label{sec:foundations}

\subsection{Set-sized data inside a class-sized field}

We work in a class theory such as NBG. The class $\No$ is a real-closed
ordered field in the class-sized sense. The usual proof that adjoining $i$
to a real-closed field gives an algebraically closed field applies to each
polynomial and yields an algebraically closed class-field
\[
 \K=\No[i].
\]
All coefficient families, supports, polynomials, and formal series appearing
in an individual construction are set-sized. A notation such as $\K[[X]]$
denotes a class of set-indexed coefficient sequences, not a set of all such
sequences. We never sum a proper-class family of nonzero terms.

Conjugation and the modulus are
\[
 \overline{x+iy}=x-iy,\qquad
 |x+iy|=(x^2+y^2)^{1/2}\in\No_{\ge0}.
\]
They obey $|zw|=|z||w|$ and the triangle inequality. For the latter, the
identity
\[
 (xu+yv)^2+(xv-yu)^2=(x^2+y^2)(u^2+v^2)
\]
gives Cauchy--Schwarz over the real-closed ordered field, and the standard
squaring argument gives $|z+w|\le |z|+|w|$.

\subsection{Normal forms and the natural valuation}

Conway normal forms, applied to the real and imaginary parts and combined,
give a unique representation
\begin{equation}\label{eq:normal}
 z=\sum_{\gamma\in S}c_\gamma t^\gamma,
 \qquad c_\gamma\in\C^*,\quad S\subset\No\text{ a well-ordered set}.
\end{equation}
The zero series has empty support. The notation $t^\gamma=\omega^{-\gamma}$
uses Conway's monomial map; it is not a claim that this is Gonshor
exponentiation with base $\omega^{-1}$.

For $z\ne0$ put $v(z)=\min S$, and put $v(0)=+\infty$.
Then
\begin{equation}\label{eq:valuation}
 v(zw)=v(z)+v(w),\qquad
 v(z+w)\ge\min\{v(z),v(w)\},\qquad v(|z|)=v(z).
\end{equation}
The last identity follows because if $z=c t^\alpha(1+\eta)$ with
$v(\eta)>0$, then $|z|=|c|t^\alpha(1+\epsilon)$ with a real infinitesimal
$\epsilon$. In particular,
\[
 z\prec w\quad\Longleftrightarrow\quad
 |z|<|w|/n\ \text{for every positive }n\in\N
 \quad\Longleftrightarrow\quad v(z)>v(w)
\]
when $w\ne0$. We use $z\prec1$ to mean that $z$ is infinitesimal.

Set
\[
 \OO_\K=\{z:v(z)\ge0\},\qquad
 \mm=\{z:v(z)>0\}.
\]
These are the valuation ring of finite elements and its maximal ideal.
Every $z\in\OO_\K$ has a unique decomposition $z=c+\epsilon$ with
$c\in\C$ and $\epsilon\in\mm$. Thus
\begin{equation}\label{eq:standardpart}
 \st:\OO_\K\longrightarrow\C,\qquad \st(c+\epsilon)=c
\end{equation}
is a surjective ring homomorphism with kernel $\mm$.

\subsection{Neighborhood limits are not sequential limits}

The neighborhood topology has balls
$B(a,r)=\{z:|z-a|<r\}$, with $r\in\No_{>0}$.
Equivalently it is generated by valuation balls $v(z-a)>\alpha$.
For example $t^{v(r)+1}\prec r$, which gives the required comparison in
one direction; a still smaller monomial ball gives the other direction.

The following elementary consequence of surreal interpolation is decisive.

\begin{proposition}[Set-sized discreteness]\label{prop:discrete}
Every set $A\subset\K$ is closed and discrete in the full neighborhood
topology. Indeed, if $A$ has at least two elements, there is a single
$r>0$ such that $|a-b|>r$ for every two distinct $a,b\in A$.
Every convergent net indexed by a set is eventually equal to its limit.
\end{proposition}
\begin{proof}
For any set $D\subset\No_{>0}$, the surreal cut $\{0\mid D\}$ supplies
an $r$ with $0<r<d$ for every $d\in D$. Apply this to the set of nonzero
distances between members of $A$. To see closedness, fix $x\notin A$ and
apply the same observation to $\{|x-a|:a\in A\}$.

For a net $(x_j)_{j\in J}$ with $J$ a set and proposed limit $x$, form
the set of its nonzero distances from $x$. A positive lower bound smaller
than all those distances makes a sufficiently small ball about $x$
contain no net values except $x$ itself. Convergence therefore implies
eventual equality. The empty-distance case is immediate.
\end{proof}

In particular, a nonterminating power series is not to be interpreted as a
limit of its finite partial sums in this topology. This does not make an
$\epsilon$--$\delta$ derivative meaningless: its quantifiers range over
all sufficiently small field elements, not over a chosen sequence.

Valuation balls are clopen. If $b$ is outside a ball, sufficiently small
changes to $b$ preserve the valuation of its difference from the center.
Consequently distinct points can be separated by clopen classes. A map
from the ordinary connected interval $[0,1]\subset\R$ to $\K$ that is
continuous for this topology must be constant. In particular an ordinary
complex circle, included in $\K$, is \emph{not} a continuous path in this
sense. Our later contours retain the ordinary complex parameter topology
and an explicitly defined coefficientwise integral.

\begin{example}[A failure of unrestricted global uniqueness]
The finite elements form a clopen class. Hence
\begin{equation}\label{eq:locallyconstant}
 L(z)=\begin{cases}0,&z\in\OO_\K,\\1,&z\notin\OO_\K\end{cases}
\end{equation}
is locally constant, everywhere analytic in the local sense developed
below, bounded, and nonconstant. Its derivative is zero everywhere. Thus
local analyticity alone cannot yield a global identity theorem, a global
``zero derivative implies constant'' theorem, or Liouville's theorem.
\end{example}

\section{Normal summation and realization of every formal germ}
\label{sec:sums}

\subsection{The summability rule}

\begin{definition}[Normal summability]\label{def:summable}
A set-indexed family $(z_j)_{j\in J}$ is \emph{normally summable} if
$\bigcup_j\supp z_j$ is well ordered and each exponent belongs to the
supports of only finitely many $z_j$. Its sum is obtained by adding the
finitely many complex coefficients at each exponent.
\end{definition}

This definition is invariant under reindexing. Products of two summable
families and regroupings of a jointly summable family have the expected
values. Joint summability, rather than a bare assertion that two iterated
sums exist, is what licenses an interchange of sums.

We use the following standard support theorem \cite{Neumann}.

\begin{lemma}[Hahn--Neumann support lemma]\label{lem:neumann}
If $A,B$ are well-ordered subsets of an ordered abelian group, then $A+B$
is well ordered and each element of $A+B$ has only finitely many
representations $a+b$ with $a\in A,b\in B$.
If $B\subset G_{>0}$ is well ordered, its additive monoid
\[
 B^*=\{0\}\cup\bigcup_{n\ge1}(B+\cdots+B)
\]
is well ordered, and a fixed group element has only finitely many
representations as a finite nondecreasing list of elements of $B$.
Consequently the families indexed by the number of factors in powers of
a series of positive valuation are normally summable.
\end{lemma}

The two-set finiteness assertion can also be seen directly: infinitely
many different first coordinates have an increasing subsequence, forcing
a decreasing subsequence of second coordinates. The monoid assertion is
the positive-support part of Neumann's lemma. All its applications here
involve set-sized ordered groups generated by the displayed supports,
even when the ambient value class is $\No$.

\begin{corollary}\label{cor:complexseries}
If $\epsilon\in\mm$ and $(c_n)_{n\ge0}$ is any complex sequence, then
$\sum_{n\ge0}c_n\epsilon^n$ is normally summable.
The same is true in finitely many infinitesimal variables.
\end{corollary}
\begin{proof}
Use the positive well-ordered union of the supports of the variables and
\cref{lem:neumann}. Complex coefficients have support $\{0\}$ or are zero;
they do not change the support estimates. For a given total degree there
are only finitely many monomials in finitely many variables.
\end{proof}

There is no growth assumption on $c_n$. Thus $\sum(n!)^2\epsilon^n$ is a
legitimate normal sum. In contrast, $\sum 2^{-n}$ is not normally summable:
infinitely many terms have exponent zero. Its ordinary complex sum can of
course be used as \emph{one coefficient}. Classical summation in the
coefficient field and Hahn summation over scales are different operations.

\subsection{An explicit scale-separation argument}

\begin{lemma}[Separated degree blocks]\label{lem:blocks}
Let $a_\nu\in\K$ be a family indexed by $\nu\in\N^d$, with $d<\infty$.
There exists $\lambda>0$ such that, with $\rho=t^\lambda$, the family
\[
 \bigl(a_\nu\rho^{|\nu|}\bigr)_{\nu\in\N^d}
\]
is normally summable. More precisely, supports belonging to different
total degrees can be chosen to lie in pairwise disjoint, increasingly
ordered blocks. For every finite tuple $u\in\OO_\K^d$, the family
$(a_\nu\rho^{|\nu|}u^\nu)_\nu$ is normally summable.
\end{lemma}
\begin{proof}
The union $S=\bigcup_\nu\supp a_\nu$ is a set, though it need not be well
ordered. Choose $\Delta>0$ with $|s|<\Delta$ for all $s\in S$, using
surreal interpolation, and then choose $\lambda>2\Delta$.
For total degree $n$ the shifted supports are contained in
\[
 (n\lambda-\Delta,n\lambda+\Delta).
\]
The upper endpoint of this interval is smaller than the lower endpoint
for degree $n+1$. Within a fixed degree there are finitely many $\nu$;
the union of their shifted well-ordered supports is well ordered, and an
exponent occurs only finitely many times. The ordered union of the degree
blocks is therefore well ordered and locally finite in the required sense.

Write $u_j=c_j+\eta_j$ with $c_j\in\C$ and $\eta_j\in\mm$. Let $B$ be
the union of the nonempty supports of the $\eta_j$. Expanding each finite
product $u^\nu$ by the binomial theorem adds elements of $B^*$ to the
previously obtained support $A$. The set $A+B^*$ is well ordered. At a
fixed exponent there are finitely many pairs from $A$ and $B^*$, finitely
many positive-support decompositions, finitely many original $\nu$ for
each member of $A$, and finitely many binomial choices for a fixed $\nu$.
This proves joint summability of the full expansion.
\end{proof}

\begin{theorem}[Realization and local boundedness]\label{thm:realization}
Every $P\in\K[[X_1,\ldots,X_d]]$ has a normally summed evaluation on a
neighborhood of zero. On a sufficiently small closed polydisc its values
are bounded in surcomplex modulus by a fixed surreal number. Formal sums,
products, derivatives, and compositions with zero constant term agree
with the corresponding operations on evaluation germs.
\end{theorem}
\begin{proof}
The evaluation assertion follows from \cref{lem:blocks} with
$X_j=\rho u_j$. The proof works for all finite $u_j$ and in particular
for $|u_j|\le1$. Every exponent in the resulting sum is at least
$-\Delta$, so its modulus is smaller than $t^{-\Delta-1}$, uniformly
on that polydisc. Cancellation can only improve this lower valuation bound.

Addition and multiplication follow by expanding jointly summable families.
For composition, write an inner series with zero constant term as a sum
of its positive-degree homogeneous parts. The fully expanded composition
has only finitely many algebraic terms at each fixed total degree. To
justify collection by degree, take a set-sized additive group $G\subset\No$
containing all supports of the original coefficients and all their finite
products. Such a group is bounded in $\No$: choose $\Delta$ larger than
$|g|$ for every $g\in G$. At a sufficiently small common scale, every
expanded term of degree $n$ has coefficient support in the $n$th separated
block. The preceding argument, including the positive supports of the
finite normalized variables, proves joint summability. Thus collecting
by total degree gives exactly the formal composite. A further shrinking
ensures that the inner values lie in a realization neighborhood of the
outer series. Derivatives are identified with actual derivatives in
\cref{thm:differentiation} below.
\end{proof}

\subsection{Microholomorphic germs}

\begin{definition}\label{def:micro}
A function $f$ on an open class is \emph{microholomorphic at $a$} if on some
neighborhood it has the form
\[
 f(a+h)=\sum_{n\ge0}a_nh^n
\]
with normal summation. A \emph{microholomorphic germ} identifies two such
functions when they agree on some neighborhood of $a$.
\end{definition}

The prefix indicates locality and does not restrict the coefficients to
finite numbers. We prove shortly that different formal series define
different germs. Thus the local category is large: it includes series
which have zero radius of ordinary complex convergence and series with
infinite surcomplex coefficients.

\section{Derivatives, Cauchy--Riemann equations, and harmonic germs}
\label{sec:calculus}

\subsection{Actual neighborhood derivatives}

The derivative is defined by the usual field-valued limit:
$f'(a)=L$ means that for each $\epsilon\in\No_{>0}$ there is a
$\delta\in\No_{>0}$ such that
\[
 0<|h|<\delta\ \Longrightarrow\
 \left|\frac{f(a+h)-f(a)}h-L\right|<\epsilon.
\]
Limits in this definition are unique, by the triangle inequality.

\begin{theorem}[Analytic differentiation]\label{thm:differentiation}
A microholomorphic function is infinitely differentiable near its center,
its derivatives are microholomorphic, and
\[
 \left(\sum_{n\ge0}a_nh^n\right)'=\sum_{n\ge1}na_nh^{n-1},\qquad
 a_n=\frac{f^{(n)}(a)}{n!}.
\]
It admits the normally summed Taylor expansion about every sufficiently
nearby point. The germ-realization map $\K[[X]]\to\{\text{germs at }a\}$
is an isomorphism of differential $\K$-algebras.
\end{theorem}
\begin{proof}
At the center write
\[
 P(h)=a_0+a_1h+h^2B(h).
\]
By \cref{thm:realization}, $|B(h)|\le M$ on a small ball. The difference
quotient differs from $a_1$ by at most $M|h|$. Taking
$|h|<\min\{r,\epsilon/(M+1)\}$ proves the derivative formula at zero.

For a nearby $b$, the proposed recentered coefficients are
\begin{equation}\label{eq:recenter}
 A_j(b)=\sum_{n\ge j}\binom nj a_n b^{n-j}.
\end{equation}
These sums exist at a sufficiently small separated scale. Expand
$a_n(b+h)^n$ by the finite binomial theorem. At total degree $n$ there
are only $n+1$ terms, so the degree-block proof and Neumann's lemma
justify the joint sum and its regrouping into
\[
 P(b+h)=\sum_{j\ge0}A_j(b)h^j.
\]
Applying the first argument at $b$ gives
$P'(b)=A_1(b)=\sum_{n\ge1}na_nb^{n-1}$.
The formal derivative is itself a realizable series. Repetition proves
all higher derivatives and the formula for the coefficients at the center.
That formula proves injectivity; surjectivity is the definition of a
microholomorphic germ. The same argument, using finite total-degree
expansions, proves the multivariable statement for mixed partial derivatives.
\end{proof}

The product and chain rules now follow either from the usual finite
algebra of difference quotients or from the formal rules and realization.
Termwise integration gives a primitive germ
\[
 \int P(X)\dd X=C+\sum_{n\ge0}\frac{a_n}{n+1}X^{n+1}.
\]
The primitive is unique up to a constant \emph{as a germ}. No global
uniqueness assertion across independent clopen pieces is implicit.

\begin{remark}[Two different derivations]
The operator $D_z$ in this article differentiates a function of a variable
and fixes every coefficient in $\K$. In particular $D_z\omega=0$ when
$\omega$ is a constant function. It is not the Berarducci--Mantova field
derivation on surreal numbers viewed as transseries; the latter has a
normalization with $\partial\omega=1$ \cite{BMderivations}. Confusing these
two roles of a surreal number would invalidate the calculus.
\end{remark}

\subsection{Cauchy--Riemann equations in the correct category}

Real microanalyticity means normal evaluation of a series in finitely many
real surreal variables. Such a function has a Fr\'echet derivative: after
subtracting its degree-zero and degree-one parts, the remainder is bounded
by $M\|h\|^2$ on a small polydisc, using the same support estimates and a
finite sum over degree-two factors.

\begin{theorem}[Cauchy--Riemann criterion]\label{thm:CR}
Let $f=u+iv$ be real microanalytic in $(x,y)$ near a point of $\No^2$.
Then $f$ is microholomorphic near that point if and only if, on a
neighborhood,
\begin{equation}\label{eq:CR}
 u_x=v_y,\qquad u_y=-v_x.
\end{equation}
At a single point these equations are exactly the condition that the
real Fr\'echet differential be complex linear.
\end{theorem}
\begin{proof}
Introduce formal coordinates $Z=X+iY$ and $W=X-iY$ over $\K$.
They are invertible linear coordinates because $2i\ne0$. Equations
\eqref{eq:CR} are equivalent to
\[
 \partial_W f=\tfrac12(\partial_X+i\partial_Y)f=0.
\]
If the equations hold as functions on a neighborhood, injectivity of
multivariable series realization makes this a formal identity. In
characteristic zero it forces every term with positive $W$-degree to
vanish. The remaining formal series depends only on $Z$, hence is
microholomorphic. Conversely a series in $Z$ is killed by $\partial_W$.
The assertion about a differential follows by applying the same calculation
to its degree-one part.
\end{proof}

This theorem does not assert that existence of two partial derivatives at
a point is sufficient for complex differentiability, or that arbitrary
pointwise complex-differentiable functions have normal-form expansions.

\begin{theorem}[Harmonic conjugate germs]\label{thm:harmonic}
The real and imaginary parts of a microholomorphic germ are harmonic:
$u_{xx}+u_{yy}=v_{xx}+v_{yy}=0$. Conversely every real microanalytic
harmonic germ in two variables has a harmonic conjugate germ, unique up
to a real constant.
\end{theorem}
\begin{proof}
Differentiate \eqref{eq:CR} and commute formal partial derivatives to
obtain the two Laplace equations. Conversely work at $(0,0)$ and set,
with formal termwise integration,
\[
 v(X,Y)=-\int_0^X u_Y(s,0)\dd s+\int_0^Y u_X(X,t)\dd t.
\]
Then $v_Y=u_X$ and, since $u_{XX}=-u_{YY}$,
\[
 v_X=-u_Y(X,0)+\int_0^Y u_{XX}(X,t)\dd t=-u_Y(X,Y).
\]
All the series involved have local realizations. Thus $u+iv$ is
microholomorphic by \cref{thm:CR}. A difference of two conjugates has
both formal partial derivatives zero and is constant.
\end{proof}

\section{Inverse functions, implicit functions, and ramification}
\label{sec:inverse}

\begin{theorem}[Inverse and implicit function theorems]\label{thm:inverse}
Let $F=(F_1,\ldots,F_d)$ be a microholomorphic germ at $a\in\K^d$.
If its Jacobian matrix $J_F(a)$ is invertible, there are open neighborhoods
on which $F$ has a microholomorphic inverse.

If $P(X,Y)$ is a vector of $q$ microholomorphic germs, $P(0,0)=0$, and
$\partial P/\partial Y(0,0)$ is an invertible $q\times q$ matrix, there
is a unique germ $Y=\varphi(X)$ with $\varphi(0)=0$ and
$P(X,\varphi(X))=0$.
\end{theorem}
\begin{proof}
Translate the centers and write $F(X)=AX+\sum_{n\ge2}F_n(X)$, with
$A$ invertible and $F_n$ homogeneous of degree $n$.
Seek $G(Y)=A^{-1}Y+\sum_{n\ge2}G_n(Y)$. Once $G_2,\ldots,G_{n-1}$
are known, the degree-$n$ part of $F(G(Y))-Y$ is $AG_n$ plus a known
homogeneous polynomial. There is exactly one choice of $G_n$ that cancels
it. This constructs a formal right inverse; the same uniqueness recursion,
or substitution into a formal left inverse, makes it two-sided.

Realize these finitely many formal series and the two composition
identities on sufficiently small neighborhoods. Analytic functions are
continuous, by the remainder estimate in the differentiation proof.
Choose a small open neighborhood $V$ of zero on which $G$ is defined,
$F\circ G=\id$, and $G(V)$ lies in a neighborhood where
$G\circ F=\id$. On the corresponding preimage neighborhood, $F$ and $G$
are inverse continuous analytic maps.

For the implicit assertion apply the inverse theorem to
$(X,Y)\mapsto(X,P(X,Y))$. Its Jacobian is block triangular with invertible
diagonal blocks. Restrict its inverse to pairs $(X,0)$. Uniqueness also
follows degree by degree from the invertible $Y$-Jacobian.
\end{proof}

The theorem is local in the full field topology. Its proof has no appeal
to completeness of a sequence of approximations. Formal coefficients are
constructed first, then normally evaluated at a sufficiently small scale.

\begin{theorem}[Finite ramification normal form]\label{thm:ramification}
Suppose $f$ is a nonconstant microholomorphic germ at $a$ and
\[
 f(a+X)-f(a)=bX^m+O(X^{m+1}),\qquad b\ne0,\quad m\ge1.
\]
There is a microholomorphic local coordinate $\phi$ with
$\phi(0)=0$, $\phi'(0)=1$, such that
\begin{equation}\label{eq:ramification}
 f(a+X)=f(a)+b\phi(X)^m.
\end{equation}
The zero at $a$ has finite order $m$ and is isolated. The germ is an open
map at $a$. On a suitable neighborhood, every sufficiently close target
different from $f(a)$ has exactly $m$ preimages, all simple.
\end{theorem}
\begin{proof}
Factor the formal series as
\[
 f(a+X)-f(a)=bX^m(1+Q(X)),\qquad Q(0)=0.
\]
The binomial series with exponent $1/m$ gives a formal unit
$(1+Q)^{1/m}$, so $\phi(X)=X(1+Q(X))^{1/m}$ has derivative one.
Realization preserves its $m$th-power identity, and \cref{thm:inverse}
makes $\phi$ a local analytic coordinate.

The factor $b(1+Q(X))$ is nonzero on a sufficiently small neighborhood,
proving isolation and order. In the $\phi$-coordinate the equation for a
target $f(a)+w$ is $u^m=w/b$. Algebraic closedness of $\K$ supplies
exactly $m$ distinct roots when $w\ne0$. Each has modulus
$(|w|/|b|)^{1/m}$, so all lie in any prescribed small coordinate
neighborhood when $|w|$ is sufficiently small. Applying the inverse
coordinate gives all preimages in the chosen neighborhood. At them the
derivative $bm u^{m-1}\phi'(X)$ is nonzero. The same root argument shows
that the image contains a neighborhood of $f(a)$.
\end{proof}

\begin{corollary}[Local identity and maximum-modulus principles]
A microholomorphic germ with infinitely high zero order is zero. If a
microholomorphic function has a local maximum of its modulus at a point,
it is constant as a germ there.
\end{corollary}
\begin{proof}
The first assertion is uniqueness of Taylor coefficients. For the second,
a nonconstant germ has an open image near the point. A neighborhood of
$b\ne0$ contains $b(1+\eta)$ with $\eta>0$, and hence a point of larger
modulus; a neighborhood of zero contains nonzero points. Either case
contradicts a local maximum.
\end{proof}

The first assertion is a germ statement, not a sequential accumulation
theorem. By \cref{prop:discrete}, no nontrivial set of points accumulates
in the full field topology. A useful global identity theorem will instead
use accumulation of \emph{standard parts}.

\section{Formal residues, local primitives, and Lagrange inversion}
\label{sec:formalres}

Let $\K((X))$ be the Laurent series with finitely many negative powers.
These are meromorphic germs: multiplication by a sufficiently large
power of $X$ makes them microholomorphic. Define
\[
 \Res_{X=0}\left(\sum_{n\ge n_0}a_nX^n\dd X\right)=a_{-1}.
\]
The residue of a meromorphic germ is therefore an element of $\K$.

\begin{proposition}[The only local primitive obstruction]\label{prop:primitive}
A Laurent differential $\alpha\in\K((X))\dd X$ has a primitive in
$\K((X))$ if and only if $\Res_0\alpha=0$. Equivalently,
\[
 \K((X))\dd X\,/\,\mathrm d\K((X))\simeq\K\,\frac{\dd X}{X}.
\]
\end{proposition}
\begin{proof}
A formal derivative has no $X^{-1}$ coefficient. Conversely, every term
$a_nX^n\dd X$ with $n\ne-1$ integrates to
$a_nX^{n+1}/(n+1)$. These coefficients form another Laurent series with
finite negative part. The only omitted term is $a_{-1}\dd X/X$.
\end{proof}

\begin{theorem}[Change of variable for residues]\label{thm:reschange}
If $g(X)=cX^m(1+O(X))$ with $c\ne0$ and $m\ge1$, then for every
$A(Y)\in\K((Y))$,
\begin{equation}\label{eq:reschange}
 \Res_{X=0}A(g(X))g'(X)\dd X
 =m\Res_{Y=0}A(Y)\dd Y.
\end{equation}
In particular residues are invariant under an invertible local coordinate.
\end{theorem}
\begin{proof}
Only finitely many negative powers of $Y$ can affect the residue, so it
suffices to check monomials. If $A(Y)=Y^n$, $n\ne-1$, the pullback is
$\mathrm d(g^{n+1})/(n+1)$, of residue zero. If $A(Y)=Y^{-1}$, write
$g=cX^mu$ with $u(0)=1$. Then
\[
 g'/g=m/X+u'/u,
\]
and $u'/u$ is a power series. Its residue is zero, giving the factor $m$.
Positive powers of $Y$ pull back to series with no negative powers and
cause no infinite contribution at exponent $-1$.
\end{proof}

\begin{corollary}[Local argument principle]
For a nonzero meromorphic germ $f=X^m u$ with $u(0)\ne0$,
\[
 \Res_0\frac{f'}f\dd X=m\in\Z.
\]
A positive $m$ is a zero order and a negative $m$ is minus a pole order.
\end{corollary}

There is also a formal Cauchy formula. In the ring
$\K((X))[[Y]]$, interpret
\[
 \frac1{X-Y}=\sum_{n\ge0}Y^nX^{-n-1}.
\]
For $P\in\K[[X]]$, coefficientwise residue in $X$ gives
\begin{equation}\label{eq:formalCauchy}
 \Res_X\frac{P(X)}{X-Y}\dd X=P(Y).
\end{equation}
The explicit ring matters: an infinite negative Laurent tail in $X$ has
not been silently declared a member of $\K((X))$.

\begin{theorem}[Lagrange--B\"urmann inversion]\label{thm:lagrange}
If $F(0)=0$, $F'(0)\ne0$, and $G=F^{-1}$ is its inverse germ, then
for $n\ge1$,
\begin{equation}\label{eq:lagrange}
 [Y^n]G(Y)=\frac1n[X^{n-1}]\left(\frac{X}{F(X)}\right)^n.
\end{equation}
More generally, for $H\in\K[[X]]$,
\[
 [Y^n]H(G(Y))=\frac1n[X^{n-1}]H'(X)
                       \left(\frac{X}{F(X)}\right)^n.
\]
These are identities of realized analytic germs as well as formal series.
\end{theorem}
\begin{proof}
For the general formula use residue change with $Y=F(X)$ to write
\[
 [Y^n]H(G(Y))=\Res_X H(X)F'(X)F(X)^{-n-1}\dd X.
\]
The residue of the derivative of $H(X)F(X)^{-n}$ vanishes. Therefore the
last expression is
\[
 \frac1n\Res_X H'(X)F(X)^{-n}\dd X
 =\frac1n[X^{n-1}]H'(X)\left(\frac X{F(X)}\right)^n.
\]
Taking $H(X)=X$ gives \eqref{eq:lagrange}. The realization and composition
theorems identify the formal inverse with the actual inverse germ.
\end{proof}

For completeness, the rational global residue theorem is already valid
over $\K$. A rational differential on the projective line has the sum of
its residues at finite poles and infinity equal to zero. Indeed, partial
fractions exist over the algebraically closed field; a term
$c/(Z-a)\dd Z$ contributes $c$ at $a$ and $-c$ at infinity, while
higher-order principal parts and polynomial differentials have zero total
residue. The analytic finite-pole theorem in \cref{sec:argument} will
reduce to this algebra after a nontrivial division step.

\section{Exponential functions, polar form, and logarithmic branches}
\label{sec:exp}

\subsection{The infinitesimal and finite exponential}

For $h\in\mm$, normal summation defines
\[
 e_{\mm}(h)=\sum_{n\ge0}\frac{h^n}{n!},\qquad
 \ell(1+h)=\sum_{n\ge1}\frac{(-1)^{n+1}}n h^n.
\]
Formal identities in finitely many infinitesimal variables, justified by
\cref{lem:neumann}, imply
\[
 e_{\mm}(h+k)=e_{\mm}(h)e_{\mm}(k),\quad
 \ell(e_{\mm}(h))=h,\quad e_{\mm}(\ell(1+h))=1+h.
\]
Thus $e_{\mm}$ is an isomorphism $(\mm,+)\to(1+\mm,\cdot)$.
For $z=c+h\in\OO_\K$, define
\[
 E_{\rm fin}(z)=e^c e_{\mm}(h).
\]
It is a group homomorphism onto $\OO_\K^*$, the finite units with nonzero
standard part, and its kernel is $2\pi i\Z$. Surjectivity follows by
writing a unit as $c(1+h)$ and choosing an ordinary logarithm of $c$;
the kernel assertion follows by first taking standard parts and then
using the injectivity of $e_{\mm}$.

\subsection{A known global exponential and its local analytic law}

Let $\exp_{\No}:\No\to\No_{>0}$ be the Kruskal--Gonshor exponential.
For a real surreal number write its normal form as
\[
 y=y_{\rm PI}+y_0+y_{\rm inf},\qquad
 \Fin(y)=y_0+y_{\rm inf}.
\]
Here the purely infinite part has positive powers of $\omega$, the
middle part is real, and the last part is infinitesimal.
The map $\Fin$ is additive. Define
\begin{equation}\label{eq:globalexp}
 \Exp(x+iy)=\exp_{\No}(x)E_{\rm fin}\bigl(i\Fin(y)\bigr).
\end{equation}
This is the canonical surcomplex exponential associated with the omnific
integer part in Ehrlich--Kaplan's treatment \cite[Section 11]{EK}.
Their existence and kernel theorem is established literature. We record
its explicit local analytic consequences.

\begin{theorem}[Global exponential and its infinitesimal expansion]
The function \eqref{eq:globalexp} is a surjective homomorphism
$(\K,+)\to(\K^*,\cdot)$, extends the ordinary complex exponential, and
satisfies
\begin{equation}\label{eq:expkernel}
 \ker\Exp=2\pi i\Oz,
 \qquad \Oz=\No_{\rm PI}+\Z.
\end{equation}
It is microholomorphic everywhere, with
\begin{equation}\label{eq:explocal}
 \Exp(a+h)=\Exp(a)\sum_{n\ge0}\frac{h^n}{n!}
 \quad(h\in\mm),\qquad \Exp'(a)=\Exp(a).
\end{equation}
\end{theorem}
\begin{proof}
Additivity of $\Fin$ and multiplicativity of the real and finite
exponentials prove the homomorphism assertion. The finite imaginary
exponential has modulus one. Hence a kernel element must have real
part zero, and its imaginary part must satisfy
$\Fin(y)\in2\pi\Z$. This is equivalent to
$y\in\No_{\rm PI}+2\pi\Z=2\pi\Oz$, since multiplication by the
nonzero real $2\pi$ preserves the class of purely infinite numbers.

For surjectivity take $z\ne0$, put $r=|z|$, and $u=z/r$.
Then $|u|=1$, so $c=\st(u)$ is an ordinary unit complex number.
Choose a real $\theta$ with $e^{i\theta}=c$ and put $v=u/c\in1+\mm$.
Conjugation commutes with the logarithm series. Since $v\bar v=1$,
\[
 \ell(v)+\overline{\ell(v)}=\ell(v\bar v)=0.
\]
Thus $\ell(v)=i\eta$ for a real infinitesimal $\eta$, and
$u=E_{\rm fin}(i(\theta+\eta))$. It follows that
$z=\Exp(\log_{\No}r+i(\theta+\eta))$.

For $h=p+iq\in\mm$ the purely infinite part of $q$ is zero.
The real exponential agrees with its Taylor series at infinitesimals.
Multiplying the two series for $p$ and $iq$, with the binomial theorem,
gives $\Exp(h)=e_{\mm}(h)$. The group law gives
\eqref{eq:explocal}, and the differentiation theorem gives its derivative.
\end{proof}

In particular $\Exp(i\omega)=1$. This is not inferred by summing
$\sum(i\omega)^n/n!$, which is not normally summable. It comes from the
chosen treatment of infinite imaginary phases.

\begin{proposition}[Local analytic axioms do not determine infinite phase]
For each real $a$, let $c_\omega(y)$ be the real coefficient of $\omega$
in the normal form of $y$, and define
\[
 E_a(x+iy)=\exp_{\No}(x)
 E_{\rm fin}\bigl(i(\Fin(y)+a\,c_\omega(y))\bigr).
\]
Each $E_a$ is an everywhere microholomorphic exponential homomorphism
extending the ordinary complex and the surreal real exponentials, with
$E_a'=E_a$, while $E_a(i\omega)=e^{ia}$.
\end{proposition}
\begin{proof}
Coefficient extraction and $\Fin$ are additive, so the group law holds.
For infinitesimal increments $q$, $c_\omega(q)=0$ and $\Fin(q)=q$.
Therefore the exact local expansion is again \eqref{eq:explocal}.
The values at $i\omega$ follow directly. On ordinary complex arguments
the extracted coefficient is zero.
\end{proof}

We use the canonical choice \eqref{eq:globalexp} in what follows.
The proposition concerns uniqueness under the listed analytic axioms;
it does not contradict the additional initial-structure normalization
used to select the canonical exponential.

\subsection{Polar form and the principal logarithm}

The proof of surjectivity also proves a finite-angle polar theorem.
Every $z\ne0$ has a representation
\[
 z=|z|E_{\rm fin}(i\theta),\qquad\theta\in\No\text{ finite},
\]
unique modulo $2\pi\Z$ among finite angles. By shifting by an ordinary
integer multiple of $2\pi$, including the infinitesimal choice at the
endpoints, one obtains a unique angle in $(-\pi,\pi]$.
For a negative real number this angle is $\pi$.

\begin{theorem}[Principal logarithm]\label{thm:log}
On $\K\setminus\No_{\le0}$ define
\[
 \Log z=\log_{\No}|z|+i\Arg z,\qquad -\pi<\Arg z<\pi.
\]
It is microholomorphic and satisfies
\[
 \Exp(\Log z)=z,\qquad \Log'(z)=1/z.
\]
The restriction of $\Exp$ to the strip
$\{x+iy:x\in\No,-\pi<y<\pi\}$ is an analytic bijection onto this
cut plane, with inverse $\Log$.
\end{theorem}
\begin{proof}
Finite-angle uniqueness gives the bijection and the inverse identities.
For a fixed $z$ off the cut and a sufficiently small $h$, the candidate
\[
 \Log z+\ell(1+h/z)
\]
still has imaginary part in $(-\pi,\pi)$. Its exponential is $z+h$.
Uniqueness in that strip therefore gives
\[
 \Log(z+h)=\Log z+
 \sum_{n\ge1}\frac{(-1)^{n+1}}{n}\left(\frac hz\right)^n.
\]
This proves microholomorphicity and the derivative. The necessary
shrinking is allowed even when the angle is only infinitesimally far
from an endpoint of the strip.
\end{proof}

\section{Coherent Hahn-holomorphic families}
\label{sec:coherent}

\subsection{Definition and evaluation on halos}

Let $U\subset\C$ be a nonempty connected open set. Write
\[
 U^\sharp=\{c+\epsilon:c\in U,\ \epsilon\in\mm\}
          =\st^{-1}(U)\subset\OO_\K.
\]
This is an open class in $\K$. It is the \emph{standard-part halo} of $U$.
For example, the halo of the ordinary unit disk consists of the finite
points whose standard part has modulus strictly less than one. It is not
the entire surcomplex disk $\{z:|z|<1\}$: points infinitesimally inside
a standard boundary point are missing from that halo.

\begin{definition}[Coherent Hahn-holomorphic family]\label{def:Hahn}
Let $\HH(U)$ consist of formal expressions
\begin{equation}\label{eq:Hahn}
 F(Z)=\sum_{\gamma\in S}f_\gamma(Z)t^\gamma,
 \qquad f_\gamma\in\OO(U),
\end{equation}
with set-sized well-ordered support
$S=\{\gamma:f_\gamma\text{ is not identically zero}\}$.
Addition and multiplication use Hahn coefficient rules. Define
\[
 D F=\sum_\gamma f_\gamma'(Z)t^\gamma.
\]
For $z=c+\epsilon\in U^\sharp$ put
\begin{equation}\label{eq:evaluation}
 F^\sharp(c+\epsilon)=
 \sum_{\gamma\in S}\sum_{n\ge0}
 t^\gamma\frac{f_\gamma^{(n)}(c)}{n!}\epsilon^n.
\end{equation}
The double sum is a single jointly normally summable family.
\end{definition}

Here ``coherent'' means that each coefficient is one ordinary holomorphic
function on the entire domain. It is not the technical claim that a
particular sheaf is coherent as a module in the sense of complex geometry.

\begin{theorem}[Faithful analytic realization]\label{thm:Hahneval}
Formula \eqref{eq:evaluation} is well defined. It gives an injective
$\K$-algebra homomorphism from $\HH(U)$ to functions on $U^\sharp$.
Every realized function is microholomorphic, and
\begin{equation}\label{eq:commuteD}
 (D F)^\sharp=(F^\sharp)'.
\end{equation}
If $F\ne0$, $\alpha=\min\supp F$, and $f_\alpha$ is its leading
coefficient function, then
\begin{equation}\label{eq:leadingreduction}
 \st\bigl(t^{-\alpha}F^\sharp(c+\epsilon)\bigr)=f_\alpha(c).
\end{equation}
\end{theorem}
\begin{proof}
Let $B=\supp\epsilon\subset\No_{>0}$. The supports of the summands in
\eqref{eq:evaluation} lie in $S+B^*$. Neumann's lemma gives well-ordering
and finiteness of the decompositions at each exponent, including the
possible powers of $\epsilon$. This proves joint summability. The same
argument for two families and the ordinary Leibniz identities prove that
evaluation preserves sums and products. Constants in $\K$ are represented
by constant coefficient functions, so the map is $\K$-linear.

At a standard point $c$, the value is $\sum_\gamma f_\gamma(c)t^\gamma$.
If all evaluated values vanish, uniqueness of normal forms makes every
$f_\gamma(c)$ zero for every $c\in U$, proving injectivity.

For analyticity at $a=c+\epsilon$, expand again with an infinitesimal
increment $h$. The common controlling support is
$S+(\supp\epsilon\cup\supp h)^*$. The finite binomial formula, followed
by allowed regrouping, gives
\[
 F^\sharp(a+h)=\sum_{k\ge0}\frac{(D^kF)^\sharp(a)}{k!}h^k.
\]
The coefficients come from the ordinary Taylor identities of the
$f_\gamma$. The intrinsic differentiation theorem proves
\eqref{eq:commuteD}. Finally, after multiplication by $t^{-\alpha}$,
the only possible exponent-zero contribution is $f_\alpha(c)$;
all other scale differences and all positive powers of $\epsilon$ have
positive valuation. This proves \eqref{eq:leadingreduction}.
\end{proof}

\begin{corollary}[Leading coefficients control possible zero shadows]
If $F\ne0$ and $F^\sharp(z)=0$, then
$f_\alpha(\st z)=0$, where $f_\alpha$ is its leading coefficient.
\end{corollary}

\subsection{Units, restriction, and composition}

For connected $U$, $\OO(U)$ is an integral domain. Hence the least support
exponent of a product is the sum of the least support exponents.
A nonzero $F=t^\alpha f_\alpha(1+E)$ is a unit in $\HH(U)$ if and only
if $f_\alpha$ is nowhere zero on $U$. When this holds,
\begin{equation}\label{eq:unitinverse}
 F^{-1}=t^{-\alpha}f_\alpha^{-1}\sum_{n\ge0}(-E)^n,
 \qquad \supp E\subset\No_{>0}.
\end{equation}
The series is normally summable by the support lemma. Conversely an
inverse forces the product of the leading coefficient functions to be
one, so the leading coefficient cannot vanish.

Restrictions to smaller ordinary domains are coefficientwise. To obtain
a sheaf on all open subsets of $\C$, impose the well-ordered-support
condition locally; on a disconnected set this permits different supports
on different components. On a connected set this local definition reduces
to \cref{def:Hahn}. Indeed, gluing over a set-sized open cover gives a
set of possible exponents and a global holomorphic function at each
exponent. A nonzero holomorphic function on a connected domain vanishes
identically on no nonempty open subset. Thus the global support is exactly
the support seen on any one nonempty chart, and is well ordered.

\begin{proposition}[Holomorphic pullback with infinitesimal perturbation]
Let $F\in\HH(V)$ and
$G=g_0+E\in\HH(U)$, where $g_0:U\to V$ is ordinary holomorphic and
$\supp E\subset\No_{>0}$. Then
\begin{equation}\label{eq:Hahncomposition}
 F\circ G:=\sum_{\gamma,n\ge0}
 t^\gamma\frac{f_\gamma^{(n)}(g_0(Z))}{n!}E(Z)^n
\end{equation}
is a member of $\HH(U)$ and
$(F\circ G)^\sharp=F^\sharp\circ G^\sharp$.
It obeys the chain rule.
\end{proposition}
\begin{proof}
The support is controlled by $\supp F+(\supp E)^*$; coefficients at
each exponent are finite sums of holomorphic functions. Thus
\eqref{eq:Hahncomposition} is defined. Expanding at a finite input and
using the common positive-support monoid proves the evaluation identity
by the ordinary Taylor composition formulas. The chain rule follows by
termwise differentiation, or by the identity of evaluations and injectivity.
\end{proof}

\subsection{Every intrinsic germ has a coherent small-scale chart}

The next result connects the two analytic layers; coherence is not limited
to convergent classical germs.

\begin{theorem}[Scale embedding of arbitrary formal germs]\label{thm:bridge}
For every $P(X)=\sum_{n\ge0}a_nX^n\in\K[[X]]$ there is a positive
monomial $\rho$ and a family $\mathcal P\in\HH(\C)$ such that
\begin{equation}\label{eq:bridge}
 \mathcal P^\sharp(\zeta)=P(\rho\zeta)
 \qquad(\zeta\in\OO_\K).
\end{equation}
Every coefficient function of $\mathcal P$ can be chosen to be a
polynomial; in the separated construction each nonzero coefficient is
in fact a constant times a single power of $Z$.
\end{theorem}
\begin{proof}
Choose $\rho=t^\lambda$ by \cref{lem:blocks}, so the supports of
$a_n\rho^n$ occupy disjoint increasing degree blocks. Expand
\[
 a_n\rho^n=\sum_{\gamma\in S_n}c_{n,\gamma}t^\gamma.
\]
The union $S=\bigcup_nS_n$ is well ordered, and every $\gamma$ belongs
to a unique degree block. Define the coefficient at that exponent to be
$c_{n,\gamma}Z^n$. This gives a Hahn family with entire polynomial
coefficients. Evaluation at finite $\zeta$ reproduces the normally
summable family $a_n\rho^n\zeta^n$, including its finite binomial
expansion about $\st\zeta$. This is exactly \eqref{eq:bridge}.
\end{proof}

A germ centered at an arbitrary $a\in\K$ is obtained by setting
$z=a+\rho\zeta$. Thus even a divergent formal series has a coherent
Cauchy calculus at a sufficiently small scale. What is not automatic is
compatibility between separately chosen charts on a large region.

\section{Coefficientwise contour integration, Cauchy, and Morera}
\label{sec:cauchy}

\subsection{Definition and elementary calculus}

\begin{definition}[External coefficientwise contour integral]\label{def:integral}
For $F=\sum_\gamma f_\gamma t^\gamma\in\HH(U)$ and an ordinary
piecewise $C^1$ curve $\Gamma$ in $U$, set
\begin{equation}\label{eq:integral}
 \int_\Gamma F(w)\dd w
 :=\sum_\gamma t^\gamma\int_\Gamma f_\gamma(w)\dd w.
\end{equation}
The integrals on the right are ordinary complex contour integrals.
The outer sum is a normal sum, with support contained in $\supp F$.
\end{definition}

This is a $\K$-linear functional, additive under concatenation and
sign-reversing under orientation reversal. For $F\in\HH(U)$ and a
curve from $a$ to $b$ with ordinary endpoints,
\[
 \int_\Gamma DF(w)\dd w=F^\sharp(b)-F^\sharp(a).
\]
All these assertions follow coefficientwise. A primitive may have an
infinite or infinitesimal value; no real-valued restriction on the answer
is present.

For this section let $D$ be a bounded Jordan domain with positively oriented
piecewise $C^1$ boundary $\Gamma$, and suppose $\overline D\subset U$.

\begin{theorem}[Cauchy's theorem and integral formula]\label{thm:Cauchy}
For $F\in\HH(U)$,
\[
 \oint_\Gamma F(w)\dd w=0.
\]
For every $z\in D^\sharp$ and $k\ge0$,
\begin{equation}\label{eq:Cauchy}
 (D^kF)^\sharp(z)=\frac{k!}{2\pi i}
       \oint_\Gamma\frac{F(w)}{(w-z)^{k+1}}\dd w.
\end{equation}
The integrand is interpreted as a Hahn-holomorphic family on a
neighborhood of $\Gamma$.
\end{theorem}
\begin{proof}
The first assertion is ordinary Cauchy's theorem for every $f_\gamma$.
Write $z=c+\epsilon$, $c\in D$, $\epsilon\in\mm$. On a neighborhood
of $\Gamma$, the expansion
\[
 \frac1{(w-c-\epsilon)^{k+1}}
 =\sum_{n\ge0}\binom{n+k}{k}
             \frac{\epsilon^n}{(w-c)^{n+k+1}}
\]
is a Hahn family. Its product with $F$ is jointly summable by
$\supp F+(\supp\epsilon)^*$. Applying ordinary Cauchy to each
coefficient gives
\begin{align*}
 \frac{k!}{2\pi i}\oint_\Gamma\frac{F(w)}{(w-z)^{k+1}}\dd w
 &=\sum_{\gamma,n} t^\gamma\epsilon^n
     k!\binom{n+k}{k}\frac{f_\gamma^{(n+k)}(c)}{(n+k)!}\\
 &=\sum_{\gamma,n}t^\gamma\frac{f_\gamma^{(n+k)}(c)}{n!}\epsilon^n
 =(D^kF)^\sharp(z).
\end{align*}
Only finite sums are combined at any fixed Hahn exponent. There is no
unjustified exchange of a real limit with a surreal sequence limit.
\end{proof}

\begin{theorem}[Coefficientwise Morera theorem]\label{thm:Morera}
Suppose $S$ is a well-ordered set of surreal exponents and $f_\gamma$
is continuous on an ordinary domain $U$ for each $\gamma\in S$.
Let $F(c)=\sum_\gamma f_\gamma(c)t^\gamma$ at standard points and
define its contour integrals coefficientwise. If
\[
 \oint_{\partial T}F(w)\dd w=0
\]
for every triangle $T$ whose closure lies in $U$, then every $f_\gamma$
is holomorphic. Thus $F$ belongs to $\HH(U)$ and has the analytic halo
realization \eqref{eq:evaluation}.
\end{theorem}
\begin{proof}
Uniqueness of normal forms makes the coefficient of $t^\gamma$ in every
triangle integral zero. Ordinary Morera's theorem applies to each
continuous $f_\gamma$. The realization theorem then applies to the
resulting holomorphic family.
\end{proof}

The continuity hypothesis here is continuity of the ordinary coefficient
functions. The theorem is not an assertion that an arbitrary function
continuous for the full surreal topology admits such coefficients.

\subsection{Cauchy estimates that genuinely retain every scale}

For an ordinary closed disk $|w-c|\le R$ contained in $U$, put
$M_\gamma(R)=\max_{|w-c|=R}|f_\gamma(w)|\in\R_{\ge0}$.
Cauchy's formula gives
\begin{equation}\label{eq:coeffCauchy}
 |f_\gamma^{(n)}(c)|\le \frac{n!M_\gamma(R)}{R^n}
 \qquad\text{for every }\gamma\text{ and }n.
\end{equation}
These are coefficientwise, scale-resolved estimates. They must not be
replaced without proof by an assertion that a single coarse surreal
bound controls every coefficient uniformly in the ordinary complex
variable. A counterexample to that replacement appears in
\cref{sec:globalprinciples}.

\subsection{Affine contours at arbitrary centers and scales}

Fix $a\in\K$ and $r\in\K^*$. A family $F\in\HH(U)$ defines
\[
 f(a+r\zeta)=F^\sharp(\zeta),\qquad\zeta\in U^\sharp.
\]
For the affine contour labeled by $a+r\Gamma$, define
\[
 \int_{a+r\Gamma}f(w)\dd w
 :=r\int_\Gamma F(\zeta)\dd\zeta.
\]
Differentiation introduces $r^{-1}$, while the differential introduces
$r$. Therefore \eqref{eq:Cauchy} transforms into precisely the usual
Cauchy formula in the affine $w$-coordinate. The center may be infinite and
$|r|$ may be infinitesimal or infinite. These are external affine
contours, with the stated coefficientwise meaning; their parameterization
is not asserted to be continuous into the full field topology.

Combined with \cref{thm:bridge}, this produces a Cauchy integral formula
for \emph{every} intrinsic microholomorphic germ on sufficiently small
affine halo circles. It is not merely a formula for extensions of
ordinary convergent power series.

\section{Constructive preparation, division, and infinitesimal Rouch\'e}
\label{sec:preparation}

This section supplies the main global result. The perturbation need not
be a convergent analytic function of a real or complex small parameter.
Its coefficients may grow arbitrarily fast, and its scales may form any
set-sized well-ordered positive surreal support.

\subsection{Ordinary division by the standard zero polynomial}

Let $f_0$ be holomorphic on a neighborhood of $\overline D$ and nonzero
on $\Gamma=\partial D$. It has finitely many zeros $c_1,\ldots,c_s$
in $D$, with positive multiplicities $m_1,\ldots,m_s$; write
\[
 d=\sum_{j=1}^s m_j,\qquad
 P_0(Z)=\prod_{j=1}^s(Z-c_j)^{m_j}.
\]
Choose an ordinary open neighborhood $V$ of $\overline D$ on which
$f_0$ has no other zeros. Then
\[
 u_0=f_0/P_0\in\OO(V)
\]
is nowhere zero. If there are no zeros, use $d=0$, $P_0=1$.

For every $h\in\OO(V)$ there is a unique polynomial $R_0h$ of degree
less than $d$ matching the derivatives of $h$ through order $m_j-1$
at each $c_j$. This is finite Hermite interpolation. The quotient
\[
 Q_0h=(h-R_0h)/P_0
\]
is holomorphic on $V$, since the numerator has the required zeros.
Thus
\begin{equation}\label{eq:division0}
 h=P_0Q_0h+R_0h.
\end{equation}
Both operators are complex linear and extend coefficientwise to Hahn
families without enlarging their support. When $d=0$, take $R_0=0$
and $Q_0=\id$. Uniqueness in \eqref{eq:division0} follows because a
polynomial of degree below $d$ divisible by $P_0$ is zero.

\subsection{Preparation by a homogeneous recursion}

\begin{theorem}[Hahn--Weierstrass preparation]\label{thm:preparation}
Suppose
\begin{equation}\label{eq:perturb}
 F=f_0+E\in\HH(U),\qquad \supp E\subset\No_{>0},
\end{equation}
where $\overline D\subset U$ and $f_0$ is nowhere zero on $\Gamma$.
After restricting to a neighborhood $V$ as above, there are unique
$p\in\K[Z]$ and $q\in\HH(V)$ with
\[
 \deg p<d,\qquad \supp p,\supp q\subset\No_{>0},
\]
such that
\begin{equation}\label{eq:preparation}
 F=u_0(1+q)W,\qquad W=P_0+p.
\end{equation}
Here $W$ is monic of degree $d$, and $u_0(1+q)$ is a unit in $\HH(V)$.
If $S=\supp(E/u_0)$, the correction supports are contained in
$S^*\setminus\{0\}$. For $d=0$ the polynomial factor is $W=1$.
\end{theorem}
\begin{proof}
Divide by $u_0$ and put $H=E/u_0$, so $F/u_0=P_0+H$.
We seek
\begin{equation}\label{eq:factorEq}
 H=P_0q+p+qp.
\end{equation}
Introduce a formal bookkeeping parameter $s$ and homogeneous corrections
$p(s)=\sum_{n\ge1}p_ns^n$, $q(s)=\sum_{n\ge1}q_ns^n$ to factor
$P_0+sH$. Define
\begin{align}\label{eq:prepRecursion}
 p_1&=R_0H,&q_1&=Q_0H,\\
 C_n&=\sum_{j=1}^{n-1}q_jp_{n-j},&
 p_n&=-R_0C_n,\qquad q_n=-Q_0C_n\quad(n\ge2).\notag
\end{align}
Every operation is defined in $\HH(V)$, and every $p_n$ is a polynomial
of degree less than $d$. Inductively
\[
 \supp p_n\cup\supp q_n\subset nS,
\]
where $nS$ is the set of sums of $n$ members of $S$. Since $S$ is
well ordered and positive, Neumann's lemma implies that the families
$(p_n)_{n\ge1}$ and $(q_n)_{n\ge1}$ are normally summable. Thus
\[
 p=\sum_{n\ge1}p_n,\qquad q=\sum_{n\ge1}q_n
\]
exist and have positive supports contained in $S^*$. This is a justified
normal evaluation at the bookkeeping value $s=1$, not a claim of
convergence of the numerical sequence of partial sums.

Equation \eqref{eq:division0} makes the degree-one terms satisfy
$H=P_0q_1+p_1$, and every degree $n\ge2$ satisfy
$0=P_0q_n+p_n+C_n$. Normal summation and multiplication therefore give
\eqref{eq:factorEq}, proving the factorization. The leading coefficient
of $u_0(1+q)$ is the nowhere-zero function $u_0$, so the factor is a
unit by \eqref{eq:unitinverse}.

For uniqueness, take two solutions and let $\alpha$ be the least
exponent in the union of the supports of their differences
$\Delta p,\Delta q$, assuming one is nonzero. The difference of the
nonlinear products $qp$ has no term of exponent $\alpha$: it is a sum
of a difference factor, of valuation at least $\alpha$, times a
positive-support factor. At exponent $\alpha$, the difference of
\eqref{eq:factorEq} consequently becomes
\[
 P_0[t^\alpha]\Delta q+[t^\alpha]\Delta p=0.
\]
Ordinary division uniqueness forces both coefficients to be zero,
contradicting the choice of $\alpha$. This proves uniqueness.
\end{proof}

The use of ordinary holomorphic functions in the division operators is
essential: it is what ties the different standard-point monads together.
The proof does not assume a norm on $\HH(V)$, nor does it assume that
$n\min S$ is cofinal in the surreal value class.

\begin{theorem}[Division by the perturbed zero polynomial]\label{thm:division}
Let $W=P_0+p$ be as in \cref{thm:preparation}. Every $A\in\HH(V)$ has
a unique decomposition
\begin{equation}\label{eq:perturbeddivision}
 A=WQ+R,\qquad Q\in\HH(V),\quad R\in\K[Z],\quad\deg R<d.
\end{equation}
\end{theorem}
\begin{proof}
Let $T(Q)=Q_0(pQ)$. This complex-linear operator shifts supports by
members of the positive set $\supp p$. Hence
\[
 Q=\sum_{n\ge0}(-T)^nQ_0A
\]
is normally summable: its support is controlled by
$\supp A+(\supp p)^*$ and the decompositions at each exponent are
finite. It satisfies $Q=Q_0A-Q_0(pQ)$. Put
$R=R_0(A-pQ)$ and use \eqref{eq:division0} to obtain
$A-pQ=P_0Q+R$, which is \eqref{eq:perturbeddivision}.

If $WQ+R=0$ and one term is nonzero, take the least exponent occurring
in $Q$ or $R$. The positive-support perturbation $pQ$ contributes
nothing at that exponent, leaving an ordinary relation
$P_0q+r=0$ with $\deg r<d$. Its uniqueness forces both coefficients to
vanish, a contradiction. The case $d=0$ is the immediate identity
$A=1\cdot A+0$.
\end{proof}

\subsection{Exact zero counting, including multiple roots}

\begin{theorem}[Infinitesimal Rouch\'e and zero clustering]\label{thm:rouche}
Under \cref{thm:preparation}, $F^\sharp$ has exactly $d$ zeros in
$D^\sharp$, counted with microholomorphic multiplicity. They are precisely
the roots in $\K$ of the monic polynomial $W$ in
\eqref{eq:preparation}. All are finite and their standard parts are among
$c_1,\ldots,c_s$.

More generally, if $F\in\HH(U)$ is nonzero, has least exponent $\alpha$,
and its leading coefficient $f_\alpha$ is nonzero on $\Gamma$, the
number of its zeros in $D^\sharp$ is exactly the ordinary zero count of
$f_\alpha$ in $D$. Adding any Hahn family all of whose exponents are
strictly greater than $\alpha$ preserves that count.
\end{theorem}
\begin{proof}
The coefficients of $W$ are finite and its standard-part polynomial is
$P_0$. If a root $\xi$ were infinite, then $v(\xi)<0$. The valuation
of its monic leading term $\xi^d$ would be strictly smaller than that
of every lower term, whose coefficient has nonnegative valuation.
A unique least-valued term cannot cancel. Thus every root is finite.
Reduction modulo $\mm$ now yields
\[
 0=\st(W(\xi))=P_0(\st\xi),
\]
so every root lies in $D^\sharp$.

Algebraic closedness gives exactly $d$ roots with algebraic multiplicity.
The unit factor in \eqref{eq:preparation} has nonzero value throughout
$V^\sharp$, by its leading coefficient. Thus its product with $W$ has
exactly the same zeros. Its germ is a unit at each point, so orders of
zeros are also unchanged. This proves the exact count, not just a
statement about approximate roots.

For the general case apply the preceding result to
$t^{-\alpha}F=f_\alpha+E$. Its nonleading exponents are positive.
A higher-valuation perturbation changes none of the leading coefficient
function, so it changes none of the zero count.
\end{proof}

Applying the theorem on disjoint small ordinary disks around the $c_j$
shows more specifically that exactly $m_j$ roots have standard part
$c_j$. Thus a multiple standard zero may split, but cannot gain or lose
multiplicity in its infinitesimal cluster.

\begin{corollary}[A classical-boundary version of Rouch\'e]
Suppose $F=t^\alpha(f_0+E)$ and $G=t^\alpha(g_0+H)$, with
$\supp E,\supp H\subset\No_{>0}$. If
\[
 |g_0(w)|<|f_0(w)|\qquad(w\in\Gamma),
\]
then $F^\sharp$ and $(F+G)^\sharp$ have the same number of zeros in
$D^\sharp$, with multiplicities.
\end{corollary}
\begin{proof}
Ordinary Rouch\'e gives the same zero count for $f_0$ and $f_0+g_0$,
both of which are nonzero on the boundary. The preceding theorem
identifies these two counts with the surcomplex counts.
\end{proof}

The strict boundary hypothesis has been stated on the controlling
coefficient functions. No claim is made that an arbitrary inequality at
an arbitrary class-sized surreal boundary is an equivalent hypothesis.

\section{Argument principle, residue theorem, and contour recovery of roots}
\label{sec:argument}

Throughout this section $\Gamma$ bounds a Jordan domain $D$ as before,
and every function is defined on an ordinary neighborhood of
$\overline D$. Quotients are interpreted as meromorphic germs at zeros
of their denominators and as Hahn families near $\Gamma$ when the
leading denominator coefficient has no boundary zero.

\begin{theorem}[Exact weighted argument principle]\label{thm:argument}
Let $F\in\HH(U)$ have a leading coefficient nonzero on $\Gamma$.
Let $\xi_1,\ldots,\xi_d$ be its zeros in $D^\sharp$, repeated according
to multiplicity. For every $A\in\HH(U)$,
\begin{equation}\label{eq:weightedargument}
 \frac1{2\pi i}\oint_\Gamma A(w)\frac{DF(w)}{F(w)}\dd w
   =\sum_{j=1}^d A^\sharp(\xi_j).
\end{equation}
In particular
\begin{equation}\label{eq:argument}
 \frac1{2\pi i}\oint_\Gamma\frac{DF(w)}{F(w)}\dd w=d\in\Z.
\end{equation}
There is no infinitesimal error term.
\end{theorem}
\begin{proof}
Preparation, including the harmless leading monomial, writes $F=UW$
on a neighborhood $V$ of $\overline D$, with $U$ a Hahn unit and
$W$ the zero polynomial. Therefore
\[
 DF/F=DU/U+W'/W.
\]
The first term is holomorphic in $\HH(V)$, so its product with $A$ has
zero contour integral. Algebraically,
\[
 W'/W=\sum_{j=1}^d\frac1{w-\xi_j}.
\]
Cauchy's formula \eqref{eq:Cauchy}, valid for each surcomplex point
$\xi_j\in D^\sharp$, makes the corresponding integral
$A^\sharp(\xi_j)$. Summing this finite family proves the formula.
\end{proof}

There is a second useful explanation of the exact integrality in
\eqref{eq:argument}. Near the boundary write
$F=t^\alpha f_\alpha(1+E)$ with $\supp E>0$. Then
\[
 DF/F=f_\alpha'/f_\alpha+D\ell(1+E).
\]
The second term has zero integral because $\ell(1+E)$ is a single-valued
Hahn family near the entire contour. Only the ordinary integer winding
contribution remains. Preparation is what identifies that integer with
all actual surcomplex roots, including split multiple roots.

\begin{corollary}[Recovering a zero cluster without solving for its roots]
Define the contour moments
\[
 s_k=\frac1{2\pi i}\oint_\Gamma w^k\frac{DF(w)}{F(w)}\dd w
 \quad(k\ge0).
\]
Then $s_0=d$ and $s_k=\sum_{j=1}^d\xi_j^k$. The monic zero-cluster
polynomial is recovered exactly by
\begin{equation}\label{eq:Newton}
 e_0=1,\qquad
 k e_k=\sum_{j=1}^k(-1)^{j-1}e_{k-j}s_j\quad(1\le k\le d),
\end{equation}
\[
 W(Z)=Z^d-e_1Z^{d-1}+e_2Z^{d-2}-\cdots+(-1)^de_d.
\]
\end{corollary}
\begin{proof}
The moment identity is \eqref{eq:weightedargument} for $A(Z)=Z^k$.
Newton's identities \eqref{eq:Newton} follow by comparing coefficients
in the logarithmic derivative of $\prod_j(1+\xi_jT)$, or directly by
formal multiplication. They hold over every characteristic-zero field.
Their recursion uniquely determines the elementary symmetric functions
of the roots, hence $W$.
\end{proof}

Thus two complementary algorithms are available: the preparation
recursion constructs $W$ by holomorphic division, while contour moments
construct the same $W$ through symmetric functions. Root multiplicities
are preserved in both descriptions.

\begin{theorem}[Finite-pole surcomplex residue theorem]\label{thm:residue}
Let $A,B\in\HH(U)$, with $B\ne0$ and the leading coefficient of $B$
nonzero on $\Gamma$. Let $M=A/B$. It has finitely many poles in
$D^\sharp$, and
\begin{equation}\label{eq:residuetheorem}
 \oint_\Gamma M(w)\dd w
   =2\pi i\sum_{\xi\in D^\sharp}\Res_{z=\xi}M(z)\dd z,
\end{equation}
where the sum is only over its actual poles and each residue is the
coefficient of $(z-\xi)^{-1}$ in its meromorphic germ.
\end{theorem}
\begin{proof}
Prepare $B=U_BW$ on a neighborhood $V$ of $\overline D$. The factor
$U_B$ is a Hahn unit, so $C=A/U_B\in\HH(V)$. By the perturbed division
theorem, $C=WQ+R$ with $Q\in\HH(V)$ and $R\in\K[Z]$ of degree
less than $\deg W$. Therefore
\begin{equation}\label{eq:meromorphicdecomp}
 M=Q+R/W.
\end{equation}
The first term is microholomorphic everywhere in $V^\sharp$ and has
zero contour integral. The second is rational over $\K$ and has a finite
partial-fraction expansion
\[
 R/W=\sum_{\xi}\sum_{k=1}^{m_\xi}
                 \frac{c_{\xi,k}}{(Z-\xi)^k}.
\]
All possible poles have standard part in $D$, by the zero theorem.
At each such point the displayed expansion agrees with its intrinsic
meromorphic germ; other pole terms are analytic on a sufficiently small
neighborhood, and cancellations simply remove zero principal parts.

Cauchy's formula for the constant coefficient function $1$ gives
\[
 \oint_\Gamma\frac{\dd w}{w-\xi}=2\pi i,
 \qquad
 \oint_\Gamma\frac{\dd w}{(w-\xi)^k}=0\quad(k\ge2).
\]
Consequently the integral of the partial fractions is
$2\pi i\sum_\xi c_{\xi,1}$, exactly the sum of the local residues.
This proves \eqref{eq:residuetheorem}.
\end{proof}

\begin{corollary}[Zeros minus poles]
If $A$ and $B$ have leading coefficient functions nonzero on $\Gamma$,
then
\[
 \frac1{2\pi i}\oint_\Gamma\frac{D(A/B)}{A/B}\dd w
  =N_A(D^\sharp)-N_B(D^\sharp),
\]
with multiplicities. Common zeros cancel equally from both sides, so the
right side also equals the number of actual zeros minus actual poles
of the quotient.
\end{corollary}

The hypotheses are important. The theorem concerns quotients of coherent
Hahn-holomorphic families. Arbitrary normally arranged collections of
meromorphic coefficients with independently moving or accumulating poles
are not being included without a common domain condition.

\section{Global principles and their exact limits}
\label{sec:globalprinciples}

\subsection{Identity and maximum-modulus principles under coherence}

The topology of $\K$ cannot provide an ordinary accumulation-point
hypothesis for a set of zeros. The standard complex coefficient domain
provides a different and effective substitute.

\begin{theorem}[Coherent identity principles]\label{thm:identity}
Let $U\subset\C$ be connected and $F\in\HH(U)$.
\begin{enumerate}[label=(\roman*)]
\item If $F^\sharp$ has zeros whose distinct standard parts have an
accumulation point in $U$, then $F=0$.
\item If $F^\sharp$ is zero as a microholomorphic germ at even one point
of $U^\sharp$, then $F=0$.
\end{enumerate}
Consequently two coherent families that agree as germs at one point
agree on the entire connected coefficient domain.
\end{theorem}
\begin{proof}
If $F\ne0$, let $f_\alpha$ be its first nonzero coefficient function.
Every zero of $F^\sharp$ has standard part among the zeros of
$f_\alpha$, by \eqref{eq:leadingreduction}. An accumulation point of
these distinct standard parts forces $f_\alpha=0$ by the ordinary
identity theorem, a contradiction. This proves (i).

For (ii), suppose the zero germ is at $a$ and put $c=\st(a)$.
If $f_\alpha(c)\ne0$, \eqref{eq:leadingreduction} already excludes a
zero at $a$. Otherwise choose a small ordinary disk $D$ about $c$ on
whose boundary $f_\alpha$ does not vanish. Its number of zeros is finite.
The preparation theorem applied to $t^{-\alpha}F$ says that $F^\sharp$
has only finitely many zeros in $D^\sharp$, counted with multiplicity.
A neighborhood on which its germ vanishes contains arbitrarily many
distinct points in the monad of $a$, a contradiction. Apply the result
to $F-G$ for the last assertion.
\end{proof}

\begin{theorem}[Maximum modulus]
A microholomorphic function whose modulus has a local maximum at $a$
is constant as a germ at $a$. If it is the realization of a member of
$\HH(U)$ with $U$ connected, it is constant on all of $U^\sharp$.
\end{theorem}
\begin{proof}
A nonconstant germ is locally open by the ramification normal form.
Every neighborhood of $f(a)$ contains a number of larger modulus:
use $f(a)(1+\eta)$ with a sufficiently small positive $\eta$ if
$f(a)\ne0$, and a nonzero sufficiently small number if $f(a)=0$.
Local openness contradicts a local maximum. In the coherent case apply
\cref{thm:identity}(ii) to $F-F^\sharp(a)$, a Hahn family obtained by
subtracting a constant in $\K$.
\end{proof}

This is a genuine neighborhood maximum principle, not merely a statement
about standard parts. The need for coherence only enters the passage
from a constant germ to one global constant.

\subsection{Liouville and polynomial growth: coefficientwise versions}

\begin{theorem}[Coefficientwise Liouville and polynomial growth]
Let $F=\sum f_\gamma t^\gamma\in\HH(\C)$.
If each $f_\gamma$ is bounded on $\C$, then $F$ is constant.
More generally, fix $m\in\N$ and suppose that for every $\gamma$ there
is a real $C_\gamma$ such that
\[
 |f_\gamma(Z)|\le C_\gamma(1+|Z|)^m\qquad(Z\in\C).
\]
Then $F$ is a polynomial of degree at most $m$ with coefficients in $\K$.
\end{theorem}
\begin{proof}
The first assertion is ordinary Liouville applied to every coefficient.
For the second, Cauchy's estimate on an ordinary circle of radius $R$
about zero gives, for $n>m$,
\[
 |f_\gamma^{(n)}(0)|/n!
 \le C_\gamma(1+R)^m/R^n.
\]
Letting the real number $R$ tend to infinity makes this derivative zero.
Thus $f_\gamma(Z)=\sum_{j=0}^m c_{\gamma,j}Z^j$. Each family
$\sum_\gamma c_{\gamma,j}t^\gamma$ has support contained in the original
well-ordered support. Regrouping the finitely many powers of $Z$ yields
$F(Z)=\sum_{j=0}^m a_j Z^j$ with $a_j\in\K$.
\end{proof}

\begin{example}[A bounded coherent function need not be constant]
Fix a positive real surreal infinitesimal $\epsilon$. The entire
coherent family $F(Z)=\epsilon Z$ realizes to a nonconstant function on
$\C^\sharp=\OO_\K$. For every finite $z$, $\epsilon z$ is
infinitesimal, so $|F^\sharp(z)|<1$. This does not meet the
coefficientwise boundedness hypothesis: its nonzero ordinary coefficient
functions are nonzero scalar multiples of $Z$ and are unbounded.

The domain here is the halo of the entire \emph{ordinary} plane, not all
of $\K$. For the latter domain, the locally constant finite/infinite
example in \cref{sec:foundations} gives a separate bounded nonconstant
microholomorphic function. These are two different obstructions to
unqualified Liouville statements.
\end{example}

\subsection{Removable singularities require the same distinction}

\begin{theorem}[Coefficientwise removal]
Let $U\subset\C$ be a domain, $c\in U$, and
$F=\sum f_\gamma t^\gamma\in\HH(U\setminus\{c\})$.
If each $f_\gamma$ is bounded in some ordinary punctured neighborhood
of $c$, then $F$ extends uniquely to a member of $\HH(U)$.
\end{theorem}
\begin{proof}
The ordinary removable-singularity theorem extends every $f_\gamma$.
Their support remains contained in the original well-ordered support,
so the extended coefficients define a Hahn family. The ordinary identity
theorem gives uniqueness coefficient by coefficient.
\end{proof}

A scalar bound on $F^\sharp$ is again insufficient. On the halo of
$0<|Z|<1$, the family $t/Z$, with $t=\omega^{-1}$, is infinitesimal
at every point and hence bounded by $1$. It has no coherent extension
across zero, because its coefficient of $t$ is $1/Z$. Notice that this
domain excludes the \emph{whole monad} of zero; it is not the class
obtained by removing just one point from the disk halo. Keeping these
two puncturing operations distinct prevents a false analogy.

\subsection{Why arbitrary coherent disk maps need not satisfy Schwarz}

Let $\mathbb D$ be the ordinary unit disk and $\epsilon>0$ a real
infinitesimal. The coherent function
\[
 F^\sharp(z)=(1+\epsilon)z\qquad(z\in\mathbb D^\sharp)
\]
fixes zero and maps $\mathbb D^\sharp$ into itself: its standard part
is $\st(z)$, which lies strictly inside the ordinary disk. Nevertheless
\[
 |(F^\sharp)'(0)|=1+\epsilon>1,
 \qquad |F^\sharp(z)|>|z|\quad(z\ne0).
\]
Thus the usual Schwarz bound is not a theorem for all coherent halo
self-maps. In the next section it is recovered for canonical classical
lifts, whose extra structure is essential.

\section{Canonical classical lifts and conformal geometry}
\label{sec:lifts}

\subsection{Functoriality of the lift}

For an ordinary holomorphic $f:U\to\C$, its \emph{canonical lift} is
\begin{equation}\label{eq:classicallift}
 f^\sharp(c+\epsilon)
 =\sum_{n\ge0}\frac{f^{(n)}(c)}{n!}\epsilon^n,
 \qquad c\in U,\quad\epsilon\in\mm.
\end{equation}
It is the realization of the Hahn family supported only at exponent
zero. In particular
\[
 \st(f^\sharp(z))=f(\st z),\qquad
 (f^\sharp)'=(f')^\sharp.
\]
If $f(U)\subset V$ and $g$ is holomorphic on $V$, then
\begin{equation}\label{eq:functor}
 (g\circ f)^\sharp=g^\sharp\circ f^\sharp.
\end{equation}
Indeed, at $z=c+\epsilon$ the difference
$f^\sharp(z)-f(c)$ is infinitesimal. Substituting its normally summed
Taylor series into the Taylor series of $g$ is licensed by the
positive-support lemma. Collection by degree is exactly the ordinary
formal chain rule, proving \eqref{eq:functor}. The identity map lifts
to the identity, and sums and products commute with lifting as well.

\begin{theorem}[Lifted conformal equivalence]
If $f:U\to V$ is an ordinary biholomorphism, then
$f^\sharp:U^\sharp\to V^\sharp$ is a bijection, both it and its
inverse are microholomorphic, and its derivative is nowhere zero.
Its inverse is $(f^{-1})^\sharp$.
\end{theorem}
\begin{proof}
Apply \eqref{eq:functor} to the two inverse identities. At $z=c+\epsilon$,
$(f')^\sharp(z)$ has nonzero standard part $f'(c)$; hence it cannot be
zero. This proves all assertions.
\end{proof}

\begin{corollary}[A standard-halo Riemann mapping theorem]
If $U\subsetneq\C$ is nonempty and simply connected, then $U^\sharp$
is analytically equivalent to $\mathbb D^\sharp$ by a canonical
classical lift.
\end{corollary}
\begin{proof}
Choose an ordinary Riemann map and apply the preceding theorem.
\end{proof}

The existence input in this corollary is precisely the ordinary Riemann
mapping theorem, not a new uniformization theorem for arbitrary
surcomplex open classes. In particular, it asserts nothing about a
classification of every open class of $\K$ using its highly disconnected
neighborhood topology.

\subsection{An exact Schwarz theorem for canonical lifts}

\begin{theorem}[Lifted Schwarz lemma]\label{thm:Schwarz}
Let $f:\mathbb D\to\mathbb D$ be ordinary holomorphic, with $f(0)=0$.
Then for every $z\in\mathbb D^\sharp$,
\[
 |f^\sharp(z)|\le |z|,\qquad |(f^\sharp)'(0)|\le1.
\]
Equality at one nonzero $z$, or equality in the derivative bound, holds
exactly when $f(Z)=aZ$ for an ordinary complex $a$ with $|a|=1$.
\end{theorem}
\begin{proof}
Ordinary Schwarz says either $f(Z)=aZ$ with $|a|=1$, or
$|f(c)|<|c|$ for every ordinary nonzero $c\in\mathbb D$ and
$|f'(0)|<1$. The rotation case lifts exactly to the same formula.
Consider the other case, and write $z=c+\epsilon$.

If $c\ne0$, then $|z|^2-|f^\sharp(z)|^2$ has positive real standard
part $|c|^2-|f(c)|^2$. It is therefore positive. If $c=0$ and
$\epsilon\ne0$, the Taylor expansion gives
\[
 f^\sharp(\epsilon)=\epsilon\bigl(f'(0)+\eta\bigr),
 \qquad \eta\in\mm.
\]
Since $|f'(0)|<1$, the modulus of the parenthesis is strictly less than
one, again giving a strict inequality. Finally, the lifted derivative
at zero is the ordinary number $f'(0)$. These observations prove the
bounds and their equality conditions.
\end{proof}

Ordinary reflection identities also lift exactly. For example, if
$U$ is conjugation-invariant and
$f(\overline Z)=\overline{f(Z)}$ on $U$, then conjugation of the
Taylor coefficients in \eqref{eq:classicallift} gives
\[
 f^\sharp(\overline z)=\overline{f^\sharp(z)}
 \qquad(z\in U^\sharp).
\]
This observation supplies the surcomplex continuation associated with
any classical Schwarz-reflection extension.

\subsection{Changing the scale}

An affine chart $z=a+r\zeta$, with arbitrary $a\in\K$ and
$r\in\K^*$, carries these results to $a+rU^\sharp$. For example,
a lifted conformal map between two such charts has the form
\[
 z\longmapsto b+s\,f^\sharp\!\left(\frac{z-a}{r}\right),
 \qquad r,s\ne0.
\]
Its derivative is $(s/r)(f')^\sharp((z-a)/r)$. Thus the construction
works at infinite centers and at arbitrarily small or large scales;
it is not restricted to a neighborhood of the embedded ordinary plane.
The separation theorem supplies such a chart for every intrinsic formal
germ, although the resulting coherent family need not be a canonical
classical lift.

\section{Worked examples and effective constructions}
\label{sec:examples}

\subsection{Normal summation without ordinary convergence}

Put $t=\omega^{-1}$. The geometric series is normally summable and
\[
 \sum_{n\ge0}t^n=\frac1{1-t}.
\]
Multiplication by $1-t$ proves the identity by cancellation of the
normally summable terms. But the ordinary sequence of partial sums
does not converge in the full neighborhood topology: its remainder
$t^{N+1}/(1-t)$ is larger than the positive number $t^\omega$ for
every ordinary integer $N$. There is no contradiction, because normal
summation was not defined by that sequence of partial sums.

Likewise
\[
 A(h)=\sum_{n\ge0}(n!)^2h^n\qquad(h\in\mm)
\]
is a microholomorphic germ, with its formally differentiated series as
its actual neighborhood derivative. At every nonzero ordinary complex
argument the displayed classical power series diverges, since the
ratio of successive absolute values is $(n+1)^2|h|$. This example
shows that the local theory is strictly larger than Taylor lifting of
ordinary analytic germs.

\subsection{Lagrange inversion with an infinite coefficient}

Let $\alpha\in\K$ be arbitrary and $F(X)=X+\alpha X^2$.
Its inverse germ is
\begin{equation}\label{eq:Catalanexample}
 G(Y)=\sum_{n\ge1}(-1)^{n-1}
       \frac1n\binom{2n-2}{n-1}\alpha^{n-1}Y^n
 =Y-\alpha Y^2+2\alpha^2Y^3-5\alpha^3Y^4+14\alpha^4Y^5-\cdots.
\end{equation}
Lagrange inversion gives the coefficient directly from
$[X^{n-1}](1+\alpha X)^{-n}/n$.
For $\alpha\ne0$ this also equals the germ
\[
 \frac{(1+4\alpha Y)^{1/2}-1}{2\alpha},
\]
where the square root has constant term one and is evaluated for
sufficiently small $Y$. Infinite $\alpha$ merely requires a smaller
scale; it does not obstruct the inverse theorem.

\subsection{A perturbed double zero and an escaping third root}

Consider the polynomial Hahn family
\begin{equation}\label{eq:cubicexample}
 F(Z)=Z^2-t+tZ^3.
\end{equation}
On the ordinary disk $|Z|<1/2$ the leading coefficient function is
$Z^2$, with two zeros counted with multiplicity. The zero-cluster theorem
therefore gives exactly two surcomplex zeros with standard part in that
disk, both infinitesimally close to zero, even though $F$ has degree three
over $\K$.

Preparation gives $F=UW$, where
\begin{align}
 W(Z)&=Z^2+(t^2+2t^5+\cdots)Z-(t+t^4+3t^7+\cdots),
       \label{eq:cubicW}\\
 U(Z)&=1+tZ-t^3-2t^6+\cdots.\label{eq:cubicU}
\end{align}
These coefficients can be found with only polynomial division by $Z^2$
in the recursion of \cref{thm:preparation}. The first homogeneous
pieces are
\[
 p_1=-t,\quad q_1=tZ,\quad p_2=t^2Z,\quad q_3=-t^3,
 \quad p_4=-t^4,\quad p_5=2t^5Z,\quad q_6=-2t^6,
 \quad p_7=-3t^7,
\]
with the unlisted earlier pieces zero.

Writing $s=t^{1/2}$ and $Z=sw$ changes the root equation to
$w^2(1+s^3w)=1$. Its two finite branches give
\begin{equation}\label{eq:cubicroots}
 \xi_\pm=\pm t^{1/2}-\frac12t^2
          \pm\frac58t^{7/2}-t^5+\cdots.
\end{equation}
The remaining root is infinite. Since the sum of all three roots of
\eqref{eq:cubicexample} is $-t^{-1}$, it is
\[
 \xi_\infty=-t^{-1}+t^2+2t^5+\cdots.
\]
Thus a higher-degree term of positive valuation can add a root outside
the finite halo without changing the zero count in that halo.

Let $\Gamma$ be the ordinary circle $|w|=1/2$, positively oriented.
The exact contour identities are
\[
 \frac1{2\pi i}\oint_\Gamma\frac{DF}{F}\dd w=2,
 \qquad
 \frac1{2\pi i}\oint_\Gamma w\frac{DF}{F}\dd w
 =\xi_++\xi_-=-t^2-2t^5-\cdots.
\]
Higher moments recover every coefficient of $W$ through Newton's
identities. The contour distinguishes the infinitesimal displacement
of the cluster's center even though both roots have standard part zero.

\subsection{A second scale beyond every finite power}

The polynomial
\[
 F(Z)=Z^2-t-t^\omega
\]
has the exact roots
\begin{equation}\label{eq:twoscale}
 \xi_\pm=\pm t^{1/2}
 \sum_{n\ge0}\binom{1/2}{n}t^{n(\omega-1)}.
\end{equation}
Indeed $t^\omega/t=t^{\omega-1}$ is infinitesimal, so the binomial
series is normally summable and its square is $1+t^{\omega-1}$.
Here $t^\omega$ is smaller than every positive ordinary integral power
of $t$. Neither a truncation at any finite $t$-order nor a sequence of
such truncations in the full field topology captures its role by an
ordinary limit. The support formulation includes it without changing
the proofs of preparation or zero counting.

\subsection{The displacement of a simple zero}

Let $F=f_0+H\in\HH(U)$ with $H$ supported at positive exponents,
$f_0(c)=0$, and $a=f_0'(c)\ne0$. In a sufficiently small ordinary disk
about $c$, the zero-cluster theorem gives exactly one root $\xi=c+\delta$.
Assume $H\ne0$ and write $\sigma=\min\supp H>0$; when $H=0$, the
root is simply $c$. The first two homogeneous perturbation orders are
\begin{equation}\label{eq:zerodisplacement}
 \delta=-\frac{H^\sharp(c)}a
 +\frac{H^\sharp(c)(DH)^\sharp(c)}{a^2}
 -\frac{f_0''(c)H^\sharp(c)^2}{2a^3}+R,
 \qquad v(R)\ge3\sigma.
\end{equation}
To prove this, insert $\delta=\delta_1+\delta_2+\cdots$ into
\[
 0=a\delta+\frac{f_0''(c)}2\delta^2+\cdots
       +H^\sharp(c)+(DH)^\sharp(c)\delta+\cdots,
\]
and assign degree one to $H$ and each of its derivatives. The degree-one
and degree-two equations give the displayed terms. At every higher
degree, division by $a$ uniquely solves the next coefficient. The
supports of the degree-$n$ terms are contained in $n\supp H$, so
Neumann's lemma sums the expansion and bounds the remaining valuation
by $3\sigma$. Formal substitution proves it is the root already
identified by preparation. This is a homogeneous support argument,
not an appeal to convergence of a sequence of Newton iterates.

\subsection{No analytic parameter-growth assumption is needed}

For a less algebraic example, let
\[
 F(Z)=Z^2-t+\sum_{n\ge2}(n!)^2t^n e^{nZ}.
\]
This belongs to $\HH(\C)$ because its scale support is a subset of
$\N$ and every coefficient is entire. No positive ordinary radius of
convergence in a numerical parameter replacing $t$ is required.
On any ordinary circle centered at zero the leading coefficient $Z^2$
is nonzero, so there are exactly two zeros in its halo, counted with
multiplicity. Their preparation polynomial and contour moments exist
as normally summed surcomplex coefficients, despite the factorial
growth of the ordinary coefficient functions in the scale index.

\subsection{What the accompanying computation verifies}

The accompanying \texttt{verify\_examples.py} performs exact rational
symbolic checks of the finite truncations in
\eqref{eq:Catalanexample}--\eqref{eq:cubicroots}. It verifies the
preparation recursion for \eqref{eq:cubicexample}, the cancellation of
the product error to the requested order, and the two displayed local
root expansions. Such calculations test signs, factors, and
coefficients. They do not certify the proper-class foundations,
normal-summability arguments, or unrestricted theorem statements.

\section{Synthesis and further directions}
\label{sec:directions}

\subsection{The resulting analytic package}

The local theory has the full algebra of formal characteristic-zero
analytic germs, but with actual function values and neighborhood
derivatives at sufficiently small surcomplex scales. The global Hahn
category adds ordinary holomorphic coherence to these scales. Its
constructive preparation theorem turns an infinitesimally perturbed
ordinary zero divisor into an exact monic polynomial over $\K$.
Division then connects the intrinsic Laurent residue to coefficientwise
contour integration, rather than leaving these as unrelated analogies.

The bridge is particularly useful: every formal surcomplex germ becomes
a Hahn family with polynomial coefficient functions after enough scale
separation. Consequently the contour and zero-cluster tools can be used
locally for germs that have no ordinary complex convergence radius.
Conversely, the counterexamples show that one cannot remove the chosen
global coherence conditions and expect all classical theorems to survive.

\subsection{A compatibility obstruction for field derivations}

One additional distinction is worth making explicit. Let $\partial$
be a derivation on $\K$ fixing $\C$, with $\partial\omega=1$;
complexification of the normalized Berarducci--Mantova derivation is
an example \cite{BMderivations}. The canonical surcomplex exponential
cannot satisfy
\begin{equation}\label{eq:expobstruction}
 \partial\Exp(z)=\Exp(z)\,\partial z\qquad\text{for every }z\in\K.
\end{equation}
Indeed $\Exp(i\omega)=1$, so the left side at $i\omega$ is zero,
whereas the right side is $i$. This proves the obstruction. More
generally, compatibility with a field derivation forces every element
of the kernel of an exponential homomorphism to be a differential
constant. The omnific-period kernel is incompatible with this
normalization of $\partial$.

There is no conflict with $D_z\Exp(z)=\Exp(z)$, proved earlier.
The latter differentiates a variable while holding every surcomplex
coefficient fixed. The former differentiates the surcomplex field's
own elements as transseries. These are different mathematical tasks.

\subsection{Relation to other analytic programs}

The coefficientwise contour functional is not proposed as the unique
notion of surreal integration. Costin and Ehrlich develop a much broader
integration program for surreal extensions of real functions
\cite{CostinEhrlich}. Their construction and the present external
holomorphic contour calculus address different input classes and
normalization questions.

There is also continuing work on Taylor expansions involving derivations,
composition, and transseries. Bagayoko and Mantova study formal Taylor
expansions over generalized power-series fields \cite{BagayokoMantova};
Mantova studies monotonicity and Taylor approximation for transseries
\cite{Mantova2026}. These topics should not be conflated with the simple
fact that a formal series in an independent variable has a sufficiently
small surreal realization. In particular, this article does not claim
that a composition law compatible with the Berarducci--Mantova field
derivation exists on all of $\No$.

\subsection{Precise extension problems suggested by the construction}

The following are proposals for further development, not theorems proved
here or claims that their broad subjects are previously unstudied.

\paragraph{Non-affine coherent atlases.}
The affine chart construction is completely explicit. A next task is to
specify classes of nonlinear chart transitions that preserve local Hahn
support control and make their coefficientwise contour functionals
intrinsic. One should formulate the overlap conditions before asserting
a global theory of surcomplex Riemann surfaces. Arbitrary topological
gluing permits the locally constant counterexamples and is too weak.

\paragraph{Several complex variables.}
The realization, Cauchy--Riemann, inverse, and implicit arguments already
work in finitely many variables. A relative preparation theorem, with
one distinguished polynomial variable and coherent parameter domains,
should require a coefficientwise division operator that preserves the
chosen parameter domain. The one-variable Hermite division used here
is not by itself a proof of that generalization.

\paragraph{Broader meromorphic and cohomological theories.}
The residue theorem covers quotients of coherent holomorphic families
with controlled boundary reduction. Allowing arbitrary meromorphic
coefficient families, essential singularities, or infinitely many poles
requires explicit locally finite pole and support conditions. Likewise,
cohomological statements would require an admissible-cover or
coefficient-domain formalism, not the ordinary connectedness of the
full neighborhood topology.

\paragraph{Effective support presentations.}
For computable generalized series, the preparation recursion and contour
moments suggest symbolic algorithms for infinitesimal root clusters.
Termination at a requested support cut, however, is a separate algorithmic
question. General surreal coefficients are not finite input data, and
one should not call the unrestricted construction an executable
algorithm without specifying an effective subfield and support language.

\paragraph{Conclusion.}
Surcomplex analysis is viable when its summation and global dependence
are made explicit. In the framework developed here, local power-series
calculus is unusually permissive, while global analytic statements
are controlled by coherent ordinary coefficient functions. Exact
preparation, zero stability, and residue identities then coexist with
infinite and infinitesimal scales, without relying on nonexistent
ordinary sequential completeness.

\appendix
\section{Applicability guide and proof dependencies}
\label{app:guide}

The table records the hypotheses that must accompany a theorem name.
Here ``intrinsic'' means a normally summed power-series germ, ``coherent''
means a member of $\HH(U)$ with the stated ordinary domain conditions,
and ``classical lift'' means \eqref{eq:classicallift}.

\small
\begin{longtable}{@{}p{0.23\textwidth}p{0.28\textwidth}p{0.438\textwidth}@{}}
\toprule
Principle & Applicable category & Essential qualification\\
\midrule
\endhead
Taylor calculus & Intrinsic & Small enough scale; normal sums, not
sequential convergence.\\[4pt]
Cauchy--Riemann & Real microanalytic & Equations on a neighborhood;
pointwise equations alone do not imply analyticity.\\[4pt]
Inverse and implicit & Intrinsic & Invertible derivative or appropriate
Jacobian block.\\[4pt]
Open mapping & Nonconstant intrinsic germ & A local assertion; does not
supply global continuation.\\[4pt]
Formal residues & Meromorphic intrinsic germs & Finite Laurent principal
part.\\[4pt]
Cauchy and Morera & Coherent & Integrals are coefficientwise on ordinary
contours, or transported affine charts.\\[4pt]
Preparation and Rouch\'e & Coherent & Leading coefficient holomorphic
near a closed ordinary disk or Jordan domain, nonzero on its boundary.\\[4pt]
Argument and residue & Coherent quotients & Boundary reduction nonzero;
finite zero clusters and the preparation/division theorem.\\[4pt]
Identity & Coherent, connected $U$ & Equal germs, or zero shadows with an
ordinary interior accumulation point.\\[4pt]
Maximum modulus & Intrinsic locally; coherent globally & Constant germ
locally; coherence propagates the constant.\\[4pt]
Liouville and growth & Coherent & Bounds on each ordinary coefficient;
a scalar surcomplex bound is insufficient.\\[4pt]
Removable singularity & Coherent & Coefficientwise boundedness at the
ordinary puncture.\\[4pt]
Schwarz lemma & Canonical classical lift & False for general coherent
self-maps of a disk halo.\\[4pt]
Riemann mapping & Canonical classical lift & Standard-part halos of
ordinary simply connected proper domains, not arbitrary open classes.\\
\bottomrule
\end{longtable}
\normalsize

The principal dependency chain is
\[
 \begin{gathered}
 \text{normal forms and Neumann supports}
 \ \Longrightarrow\ \text{germ realization and coherent evaluation},\\
 \text{ordinary Hermite division and positive supports}
 \ \Longrightarrow\ \text{preparation and division},\\
 \text{preparation and algebraic closure}
 \ \Longrightarrow\ \text{exact zero clusters},\\
 \text{coefficientwise Cauchy and preparation/division}
 \ \Longrightarrow\ \text{argument and residue formulas}.
 \end{gathered}
\]
Ordinary complex analysis enters through the ordinary coefficient
functions. Surreal arithmetic enters through normal summation and
algebraic closure. Keeping those two inputs separate is the central
organizational principle of the article.

\begin{thebibliography}{99}

\bibitem{Ahlfors}
Lars V. Ahlfors,
\emph{Complex Analysis}, third edition,
McGraw--Hill, 1979.

\bibitem{Alling}
Norman L. Alling,
\emph{Foundations of Analysis over Surreal Number Fields},
North-Holland Mathematics Studies, vol.~141,
North-Holland, 1987.
See also ``Fundamentals of analysis over surreal number fields,''
\emph{Rocky Mountain Journal of Mathematics} \textbf{19} (1989), 565--573,
\href{https://doi.org/10.1216/RMJ-1989-19-3-565}{doi:10.1216/RMJ-1989-19-3-565}.

\bibitem{BagayokoMantova}
Vincent Bagayoko and Vincenzo Mantova,
``Taylor expansions over generalised power series,''
preprint, 2025,
\href{https://arxiv.org/abs/2509.08473}{arXiv:2509.08473}.

\bibitem{BMderivations}
Alessandro Berarducci and Vincenzo Mantova,
``Surreal numbers, derivations and transseries,''
\emph{Journal of the European Mathematical Society}
\textbf{20} (2018), 339--390.
\href{https://arxiv.org/abs/1503.00315}{arXiv:1503.00315}.

\bibitem{BMgerms}
Alessandro Berarducci and Vincenzo Mantova,
``Transseries as germs of surreal functions,''
\emph{Transactions of the American Mathematical Society}
\textbf{371} (2019), 3549--3592.
\href{https://arxiv.org/abs/1703.01995}{arXiv:1703.01995};
in particular Section~3 on surreal analytic functions.

\bibitem{Conway}
John H. Conway,
\emph{On Numbers and Games}, second edition,
A K Peters, 2001; first edition, 1976.

\bibitem{CostinEhrlich}
Ovidiu Costin and Philip Ehrlich,
``Integration on the surreals,''
\emph{Advances in Mathematics} \textbf{452} (2024), article 109823.
\href{https://arxiv.org/abs/2208.14331}{arXiv:2208.14331}.

\bibitem{EK}
Philip Ehrlich and Elliot Kaplan,
``Surreal ordered exponential fields,''
\emph{The Journal of Symbolic Logic} \textbf{86} (2021), 1066--1115,
\href{https://doi.org/10.1017/jsl.2021.59}{doi:10.1017/jsl.2021.59}.
\href{https://arxiv.org/abs/2002.07739}{arXiv:2002.07739};
Section~11 treats surcomplex exponentiation.

\bibitem{Gonshor}
Harry Gonshor,
\emph{An Introduction to the Theory of Surreal Numbers},
London Mathematical Society Lecture Note Series, vol.~110,
Cambridge University Press, 1986.

\bibitem{Mantova2026}
Vincenzo Mantova,
``Monotonicity and a Taylor approximation theorem for transseries,''
preprint, 2026,
\href{https://arxiv.org/abs/2601.07747}{arXiv:2601.07747},
version~2, May~9, 2026.

\bibitem{Neumann}
Bernhard H. Neumann,
``On ordered division rings,''
\emph{Transactions of the American Mathematical Society}
\textbf{66} (1949), 202--252,
\href{https://doi.org/10.1090/S0002-9947-1949-0032593-5}
{doi:10.1090/S0002-9947-1949-0032593-5}.

\end{thebibliography}
\end{document}
